% Options for packages loaded elsewhere
% Options for packages loaded elsewhere
\PassOptionsToPackage{unicode}{hyperref}
\PassOptionsToPackage{hyphens}{url}
\PassOptionsToPackage{dvipsnames,svgnames,x11names}{xcolor}
%
\documentclass[
  12pt,
  letterpaper,
  12pt,
  oneside,
  english,
  onehalfspacing,
  headsepline]{MastersDoctoralThesis}
\usepackage{xcolor}
\usepackage[paper=a4paper,inner=2.5cm,outer=3.8cm,top=2cm,bottom=2cm]{geometry}
\usepackage{amsmath,amssymb}
\setcounter{secnumdepth}{5}
\usepackage{iftex}
\ifPDFTeX
  \usepackage[T1]{fontenc}
  \usepackage[utf8]{inputenc}
  \usepackage{textcomp} % provide euro and other symbols
\else % if luatex or xetex
  \usepackage{unicode-math} % this also loads fontspec
  \defaultfontfeatures{Scale=MatchLowercase}
  \defaultfontfeatures[\rmfamily]{Ligatures=TeX,Scale=1}
\fi
\usepackage{lmodern}
\ifPDFTeX\else
  % xetex/luatex font selection
\fi
% Use upquote if available, for straight quotes in verbatim environments
\IfFileExists{upquote.sty}{\usepackage{upquote}}{}
\IfFileExists{microtype.sty}{% use microtype if available
  \usepackage[]{microtype}
  \UseMicrotypeSet[protrusion]{basicmath} % disable protrusion for tt fonts
}{}
\usepackage{setspace}
\makeatletter
\@ifundefined{KOMAClassName}{% if non-KOMA class
  \IfFileExists{parskip.sty}{%
    \usepackage{parskip}
  }{% else
    \setlength{\parindent}{0pt}
    \setlength{\parskip}{6pt plus 2pt minus 1pt}}
}{% if KOMA class
  \KOMAoptions{parskip=half}}
\makeatother
% Make \paragraph and \subparagraph free-standing
\makeatletter
\ifx\paragraph\undefined\else
  \let\oldparagraph\paragraph
  \renewcommand{\paragraph}{
    \@ifstar
      \xxxParagraphStar
      \xxxParagraphNoStar
  }
  \newcommand{\xxxParagraphStar}[1]{\oldparagraph*{#1}\mbox{}}
  \newcommand{\xxxParagraphNoStar}[1]{\oldparagraph{#1}\mbox{}}
\fi
\ifx\subparagraph\undefined\else
  \let\oldsubparagraph\subparagraph
  \renewcommand{\subparagraph}{
    \@ifstar
      \xxxSubParagraphStar
      \xxxSubParagraphNoStar
  }
  \newcommand{\xxxSubParagraphStar}[1]{\oldsubparagraph*{#1}\mbox{}}
  \newcommand{\xxxSubParagraphNoStar}[1]{\oldsubparagraph{#1}\mbox{}}
\fi
\makeatother


\usepackage{longtable,booktabs,array}
\usepackage{calc} % for calculating minipage widths
% Correct order of tables after \paragraph or \subparagraph
\usepackage{etoolbox}
\makeatletter
\patchcmd\longtable{\par}{\if@noskipsec\mbox{}\fi\par}{}{}
\makeatother
% Allow footnotes in longtable head/foot
\IfFileExists{footnotehyper.sty}{\usepackage{footnotehyper}}{\usepackage{footnote}}
\makesavenoteenv{longtable}
\usepackage{graphicx}
\makeatletter
\newsavebox\pandoc@box
\newcommand*\pandocbounded[1]{% scales image to fit in text height/width
  \sbox\pandoc@box{#1}%
  \Gscale@div\@tempa{\textheight}{\dimexpr\ht\pandoc@box+\dp\pandoc@box\relax}%
  \Gscale@div\@tempb{\linewidth}{\wd\pandoc@box}%
  \ifdim\@tempb\p@<\@tempa\p@\let\@tempa\@tempb\fi% select the smaller of both
  \ifdim\@tempa\p@<\p@\scalebox{\@tempa}{\usebox\pandoc@box}%
  \else\usebox{\pandoc@box}%
  \fi%
}
% Set default figure placement to htbp
\def\fps@figure{htbp}
\makeatother


% definitions for citeproc citations
\NewDocumentCommand\citeproctext{}{}
\NewDocumentCommand\citeproc{mm}{%
  \begingroup\def\citeproctext{#2}\cite{#1}\endgroup}
\makeatletter
 % allow citations to break across lines
 \let\@cite@ofmt\@firstofone
 % avoid brackets around text for \cite:
 \def\@biblabel#1{}
 \def\@cite#1#2{{#1\if@tempswa , #2\fi}}
\makeatother
\newlength{\cslhangindent}
\setlength{\cslhangindent}{1.5em}
\newlength{\csllabelwidth}
\setlength{\csllabelwidth}{3em}
\newenvironment{CSLReferences}[2] % #1 hanging-indent, #2 entry-spacing
 {\begin{list}{}{%
  \setlength{\itemindent}{0pt}
  \setlength{\leftmargin}{0pt}
  \setlength{\parsep}{0pt}
  % turn on hanging indent if param 1 is 1
  \ifodd #1
   \setlength{\leftmargin}{\cslhangindent}
   \setlength{\itemindent}{-1\cslhangindent}
  \fi
  % set entry spacing
  \setlength{\itemsep}{#2\baselineskip}}}
 {\end{list}}
\usepackage{calc}
\newcommand{\CSLBlock}[1]{\hfill\break\parbox[t]{\linewidth}{\strut\ignorespaces#1\strut}}
\newcommand{\CSLLeftMargin}[1]{\parbox[t]{\csllabelwidth}{\strut#1\strut}}
\newcommand{\CSLRightInline}[1]{\parbox[t]{\linewidth - \csllabelwidth}{\strut#1\strut}}
\newcommand{\CSLIndent}[1]{\hspace{\cslhangindent}#1}



\setlength{\emergencystretch}{3em} % prevent overfull lines

\providecommand{\tightlist}{%
  \setlength{\itemsep}{0pt}\setlength{\parskip}{0pt}}



 


% Prevent babel option clash with MastersDoctoralThesis.cls
\PassOptionsToPackage{english}{babel}

% preamble.tex — LaTeX packages and thesis macros
\usepackage{graphicx}
\usepackage[T1]{fontenc}
\usepackage[utf8]{inputenc}
\usepackage[sc]{mathpazo}
\linespread{1.05}

\usepackage{xcolor}
\usepackage{hyperref}
\usepackage{setspace}
\usepackage{booktabs}
\usepackage{longtable}
\usepackage{enumitem}
\usepackage{amsmath}
\usepackage{amssymb}

% cosupervisor is not defined in MastersDoctoralThesis.cls — define manually
\NewDocumentCommand{\cosupervisor}{m}{\newcommand{\cosupname}{#1}}

% keywordnames is only created when \keywords{} is called — provide safe fallback
\providecommand{\keywordnames}{}

% Section numbering depth
\setcounter{secnumdepth}{3}
\setcounter{tocdepth}{3}

% Redefine chapter heading: "1  Title" on a single line, no "Chapter" prefix
% Must use \makeatletter outside \AtBeginDocument to avoid hook expansion issues
\makeatletter
\renewcommand{\@makechapterhead}[1]{%
  \vspace*{20pt}%
  {\parindent 0pt \raggedright \normalfont
    \ifnum \c@secnumdepth >\m@ne
      \if@mainmatter
        \Huge\bfseries \thechapter \quad #1
      \else
        \Huge\bfseries #1
      \fi
    \else
      \Huge\bfseries #1
    \fi
    \par\nobreak
    \vspace*{40pt}%
  }%
  \thispagestyle{plain}%
}
\makeatother

\AtBeginDocument{
  % PDF metadata and link colours
  \hypersetup{
    pdftitle={\ttitle},
    pdfauthor={\authorname},
    allcolors=black,
    urlcolor=blue
  }

  % Number headings from chapter level: 1 / 1.1 / 1.1.1 / 1.1.1.1
  \renewcommand{\thechapter}{\arabic{chapter}}
  \renewcommand{\thesection}{\thechapter.\arabic{section}}
  \renewcommand{\thesubsection}{\thesection.\arabic{subsection}}
  \renewcommand{\thesubsubsection}{\thesubsection.\arabic{subsubsection}}
}
\makeatletter
\@ifpackageloaded{bookmark}{}{\usepackage{bookmark}}
\makeatother
\makeatletter
\@ifpackageloaded{caption}{}{\usepackage{caption}}
\AtBeginDocument{%
\ifdefined\contentsname
  \renewcommand*\contentsname{Table of contents}
\else
  \newcommand\contentsname{Table of contents}
\fi
\ifdefined\listfigurename
  \renewcommand*\listfigurename{List of Figures}
\else
  \newcommand\listfigurename{List of Figures}
\fi
\ifdefined\listtablename
  \renewcommand*\listtablename{List of Tables}
\else
  \newcommand\listtablename{List of Tables}
\fi
\ifdefined\figurename
  \renewcommand*\figurename{Figure}
\else
  \newcommand\figurename{Figure}
\fi
\ifdefined\tablename
  \renewcommand*\tablename{Table}
\else
  \newcommand\tablename{Table}
\fi
}
\@ifpackageloaded{float}{}{\usepackage{float}}
\floatstyle{ruled}
\@ifundefined{c@chapter}{\newfloat{codelisting}{h}{lop}}{\newfloat{codelisting}{h}{lop}[chapter]}
\floatname{codelisting}{Listing}
\newcommand*\listoflistings{\listof{codelisting}{List of Listings}}
\makeatother
\makeatletter
\makeatother
\makeatletter
\@ifpackageloaded{caption}{}{\usepackage{caption}}
\@ifpackageloaded{subcaption}{}{\usepackage{subcaption}}
\makeatother
\usepackage{bookmark}
\IfFileExists{xurl.sty}{\usepackage{xurl}}{} % add URL line breaks if available
\urlstyle{same}
\hypersetup{
  pdftitle={Measuring Technical Debt in Mission Critical Trading Systems},
  pdfauthor={Tobias Zeier},
  colorlinks=true,
  linkcolor={black},
  filecolor={Maroon},
  citecolor={black},
  urlcolor={blue},
  pdfcreator={LaTeX via pandoc}}


\thesistitle{Measuring Technical Debt in Mission Critical Trading
Systems}
\supervisor{Dr Zahid Ullah}
\cosupervisor{Doug Millward}
\degree{Master of Science}
\author{Tobias Zeier}
\university{University of Essex Online}
\department{Computing Department}
\group{Enterprise IT Management Programme}
\faculty{}

\setcounter{tocdepth}{1}
\begin{document}
\frontmatter % Use roman page numbering style (i, ii, iii, iv...) for the pre-content pages

\pagestyle{plain} % Default to the plain heading style until the thesis style is called for the body content

%----------------------------------------------------------------------------------------
%	TITLE PAGE
%----------------------------------------------------------------------------------------

\begin{titlepage}
\begin{center}

\vspace*{.06\textheight}
{\scshape\LARGE \univname\par}\vspace{1.5cm} % University name
\textsc{\Large \degreename\ Thesis}\\[0.5cm] % Thesis type

\HRule \\[0.4cm] % Horizontal line
{\LARGE \bfseries \ttitle\par}\vspace{0.2cm} % Thesis title
{\normalsize\bfseries A Service Level Quantitative
Framework\par}\vspace{0.2cm} %Thesis subtitle
\HRule \\[1.5cm] % Horizontal line
 
\begin{minipage}[t]{0.4\textwidth}
\begin{flushleft} \normalsize
\emph{Author:}\\
\authorname
\end{flushleft}
\end{minipage}
\begin{minipage}[t]{0.4\textwidth}
\begin{flushright} \normalsize
\emph{Supervisor:} \\
%
\supname\\[0.3cm]

\emph{Co-Supervisor:}\\
\cosupname
\end{flushright}
\end{minipage}\\[3cm]
 
\vfill

\normalsize \textit{A thesis submitted in fulfillment of the requirements\\ for the degree of \degreename}\\[0.3cm] % University requirement text
\textit{in the}\\[0.4cm]
\groupname\\
\deptname\\[1cm] % Research group name and department name
 
\vfill


\includegraphics[height=2.5cm]{images/logo.png} % University/department logo

{\normalsize 21 April 2026}\\[1cm] % Date

 
\vfill
\end{center}
\end{titlepage}

%----------------------------------------------------------------------------------------
%	DECLARATION PAGE
%----------------------------------------------------------------------------------------
\begin{declaration}
\addchaptertocentry{\authorshipname} % Add the declaration to the table of contents
\noindent I, \authorname, declare that this thesis titled \textit{\ttitle}, and the work presented in it are my own. I confirm that:

\begin{itemize}
  \item This work was done wholly during the course of the computing project module as part of my MSc Enterprise IT Management degree at the University of Essex Online.
  \item Where any part of this thesis has previously been submitted for a degree or any other qualification at this University or any other institution, this has been clearly stated.
  \item Where I have consulted the published work of others, this is always clearly attributed.
  \item Where I have quoted from the work of others, the source is always given. With the exception of such quotations, this thesis is entirely my own work.
  \item I have acknowledged all main sources of help.
  \item I have made use of AI tools in the preparation of this thesis. The tools employed and the nature of their use are detailed below:
    \begin{itemize}
      \item \textbf{Tools used:} Claude AI, Google Gemini, and Grammarly.
      \item \textbf{Purpose:} I have used these tools solely to improve the linguistic quality and readability of this thesis, including the correction of grammatical errors, the refinement of sentence structures, and the maintenance of a consistent academic tone throughout. All core concepts, research data, analysis, and conclusions are my own.
      \item \textbf{Scope:} I have applied AI assistance across the full thesis for grammar checking and in selected sections to improve sentence structure and clarity.
    \end{itemize}
  \item I acknowledge full responsibility for the entire content of this thesis, including those sections where AI assistance was employed, and accept accountability for any violations of ethical standards in publications.
\end{itemize}

\vspace{1cm}

\noindent Signed:\\
\rule[0.5em]{25em}{0.5pt}

\noindent Date:\\
\rule[0.5em]{25em}{0.5pt}


\end{declaration}

\clearpage


%----------------------------------------------------------------------------------------
%	ABSTRACT PAGE
%----------------------------------------------------------------------------------------

\begin{abstract}
\addchaptertocentry{\abstractname} % Add the abstract to the table of contents
This thesis develops a service-level quantitative framework for
measuring and prioritising technical debt in mission-critical trading
systems. Existing measurement approaches rely predominantly on static
code analysis, which is inapplicable to vendor-dominated enterprise
environments where operational risks manifest in service instability,
patch delay and infrastructure obsolescence rather than source code. The
framework applies TOPSIS to four operational indicators, incident
frequency, mean time to restore, change failure rate and patch recency,
derived from live ITSM and CMDB records at a Swiss securities trading
institution. Indicator selection is grounded in a critical review of
technical debt theory, measurement frameworks, financial enterprise IT
governance, and multi-criteria decision-making. The study adopts a
Design Science Research paradigm and a pragmatist stance, treating the
relationship between degraded operational indicators and technical debt
severity as a testable hypothesis rather than an established premise.
The framework's evaluation centres on ranking stability under criterion
weight perturbation, contributing a reproducible prioritisation
instrument that bridges IT service management and software quality
engineering in a domain currently underserved by the technical debt
literature.
\end{abstract}


%----------------------------------------------------------------------------------------
%	ACKNOWLEDGEMENTS
%----------------------------------------------------------------------------------------

\begin{acknowledgements}
\addchaptertocentry{\acknowledgementname} % Add the acknowledgements to the table of contents
\vspace{1cm}
I would like to express my sincere gratitude to Dr Zahid Ullah for his expert guidance, patience, and continued support throughout this research project. His insightful feedback and encouragement were invaluable in shaping this work.

I am equally grateful to Doug Millward for his thoughtful co-supervision during the early stages of this thesis and his practical perspectives on the subject matter.

I would also like to thank the University of Essex Online for providing an excellent academic environment and the resources necessary to complete this degree.

Finally, I extend my heartfelt thanks to my colleagues, friends, and family for their understanding and encouragement throughout this journey.

\end{acknowledgements}


% 

\begingroup
\hypersetup{linkcolor=black}

\tableofcontents % Prints the main table of contents

\listoffigures % Prints the list of figures

\listoftables % Prints the list of tables

\endgroup


%----------------------------------------------------------------------------------------
%	ABBREVIATIONS
%----------------------------------------------------------------------------------------

\clearpage
\begin{abbreviations}{ll}
\textbf{AHP}    & Analytic Hierarchy Process \\
\textbf{BCS}    & British Computer Society \\
\textbf{CFR}    & Change Failure Rate \\
\textbf{CI/CD}  & Continuous Integration / Continuous Delivery \\
\textbf{CMDB}   & Configuration Management Database \\
\textbf{COTS}   & Commercial Off the Shelf \\
\textbf{DORA}   & DevOps Research and Assessment \\
\textbf{DSR}    & Design Science Research \\
\textbf{FINMA}  & Swiss Financial Market Supervisory Authority \\
\textbf{FADP}   & Federal Act on Data Protection \\
\textbf{IT}     & Information Technology \\
\textbf{ITIL}   & IT Infrastructure Library \\
\textbf{ITSM}   & IT Service Management \\
\textbf{MCDM}   & Multi-Criteria Decision-Making \\
\textbf{MTTR}   & Mean Time to Restore \\
\textbf{P1}     & Priority 1 \\
\textbf{P2}     & Priority 2 \\
\textbf{P3}     & Priority 3 \\
\textbf{RQ}     & Research Question \\
\textbf{SLA}    & Service Level Agreement \\
\textbf{SQALE}  & Software Quality Assessment based on Lifecycle Expectations \\
\textbf{TD}     & Technical Debt \\
\textbf{TOPSIS} & Technique for Order of Preference by Similarity to Ideal Solution \\
\textbf{VIKOR}  & VlseKriterijumska Optimizacija I Kompromisno Resenje \\
\end{abbreviations}



%----------------------------------------------------------------------------------------
%	SYMBOLS
%----------------------------------------------------------------------------------------

\clearpage
\begin{symbols}{ll}
$C_i$              & TOPSIS closeness coefficient \\
$S_i^{\text{+}}$   & Euclidean distance from positive ideal solution \\
$S_i^{\text{-}}$   & Euclidean distance from negative ideal solution \\
$x_{ij}$           & Raw indicator value for service ${i}$, criterion ${j}$ \\
$r_{ij}$           & Vector-normalised indicator value \\
$v_{ij}$           & Weighted normalised indicator value \\
$w_j$              & Criterion weight \\
$m$                & Number of services evaluated \\
$n$                & Number of indicators \\
\end{symbols}




%----------------------------------------------------------------------------------------
%	THESIS CONTENT - CHAPTERS
%----------------------------------------------------------------------------------------

\mainmatter % Begin numeric (1,2,3...) page numbering

\pagestyle{thesis} % Return the page headers back to the "thesis" style
\newcommand{\keyword}[1]{\textbf{#1}}
\newcommand{\tabhead}[1]{\textbf{#1}}
\newcommand{\code}[1]{\texttt{#1}}
\newcommand{\file}[1]{\texttt{\bfseries#1}}
\newcommand{\option}[1]{\texttt{\itshape#1}}


\setstretch{1.5}
\bookmarksetup{startatroot}

\chapter{Introduction}\label{sec-chapter1}

This chapter establishes the research context, defines the problem the
study addresses, and sets out the aims, objectives, research questions,
and significance of the work. It follows the structure of the thesis as
a whole, moving from the broad territory of technical debt in enterprise
IT to the specific gap in knowledge that the proposed framework is
designed to fill.

\section{Background and Motivation}\label{background-and-motivation}

Technical debt describes the implied future cost of accepting short-term
technical compromises in place of more robust solutions (Cunningham,
1992). Since its introduction, the concept has expanded from source code
quality to encompass architectural, design, test and documentation
dimensions of complex IT systems (Kruchten, Nord and Ozkaya, 2012). Its
operational and financial consequences are measurable. Banker, Liang and
Ramasubbu (2020) found that technical debt significantly reduced firm
performance over an eleven-year product lifecycle. The penalty more than
doubled as debt compounded, and firms with median debt levels saw their
system's net value turn negative by year five. In financial trading
environments, these consequences are intensified by continuous
availability requirements, vendor-dominated architectures and regulatory
oversight. This makes the absence of a quantitative prioritisation
mechanism particularly costly (Ramasubbu and Kemerer, 2021; Haki
\emph{et al.}, 2023).

\section{Problem Statement}\label{problem-statement}

The technical debt literature lacks a standardised, quantitative
framework for measuring and prioritising debt at the service level in
vendor-dominated enterprise IT. Lenarduzzi \emph{et al.} (2021) reviewed
44 primary studies on prioritisation and found that quantitative
multi-criteria approaches are uncommon and empirical validation on
industrial datasets is rare. Dominant measurement tools operationalise
debt as remediation effort derived from source code rule violations
(Biazotto \emph{et al.}, 2024), rendering them inapplicable where source
code access is unavailable, as is typical in trading IT. Even where code
is accessible, tools applied to the same codebase produce substantially
divergent estimates (Amanatidis \emph{et al.}, 2020), leaving
infrastructure lifecycle debt, platform obsolescence and vulnerability
debt largely unaddressed. Although DORA metrics have been validated as
predictors of software delivery performance (Forsgren, Humble and Kim,
2018), the inferential step from degraded operational indicators to
elevated technical debt severity is not yet empirically established and
is treated in this study as a testable hypothesis. To the best of
current knowledge, no existing framework applies multi-criteria
decision-making to live operational service management data in order to
produce a ranked technical debt prioritisation instrument for financial
trading IT. This gap defines the research problem.

\section{Research Aims and
Objectives}\label{research-aims-and-objectives}

This study develops a service-level quantitative framework for measuring
and prioritising technical debt in mission-critical trading systems,
grounded in operational service management data rather than source code
analysis. The specific objectives are as follows.

\begin{enumerate}
\def\labelenumi{\arabic{enumi}.}
\tightlist
\item
  To review literature on technical debt measurement, prioritisation,
  financial enterprise IT governance and multi-criteria decision-making.
\item
  To identify and operationally define quantifiable technical debt
  indicators derived from live ITSM and CMDB records.
\item
  To design and apply a TOPSIS-based multi-criteria model for scoring
  and ranking services by technical debt priority.
\item
  To evaluate the framework assessing ranking stability under criterion
  weight perturbation.
\end{enumerate}

\section{Research Questions}\label{research-questions}

The study is guided by the following primary research question.

\begin{quote}
RQ1: \emph{How can technical debt in mission-critical trading systems be
measured quantitatively at the service level?}
\end{quote}

Two secondary questions support the investigation.

\begin{itemize}
\tightlist
\item
  RQ2: To what extent do operational service management metrics serve as
  valid proxies for technical debt severity in vendor-dominated
  enterprise IT environments?
\item
  RQ3: How stable are the resulting service rankings when criterion
  weights are perturbed?
\end{itemize}

\section{Significance of the Study}\label{significance-of-the-study}

The principal academic contribution is a prescriptive design artefact
that bridges IT service management and software quality engineering in a
domain currently underserved by the technical debt literature. To the
best of current knowledge, no prior study has applied TOPSIS to live
operational records from a financial trading institution for technical
debt prioritisation. Practically, the framework operates on data
routinely available in enterprise ITSM and CMDB systems without
requiring source code access, making it directly applicable to
vendor-dominated trading IT architectures and supporting auditable
remediation decisions in regulated environments (\emph{{FINMA Circular}
2023/1 - {Operational} risks and resilience -- banks}, 2024).

\section{Thesis Structure}\label{thesis-structure}

Chapter~\ref{sec-chapter2} reviews the theoretical and empirical
foundations across four bodies of literature and identifies the research
gap. Chapter 3 describes the research design, methodology and ethical
considerations. Chapter~\ref{sec-chapter4} presents the framework
development and results, including TOPSIS computation and sensitivity
analysis. Chapter~\ref{sec-chapter5} evaluates the framework against the
criteria established in Chapter~\ref{sec-chapter3}.
Chapter~\ref{sec-chapter6} concludes with responses to the research
questions, a critical self-evaluation, the academic and practical
contributions, and directions for future work.

\bookmarksetup{startatroot}

\chapter{Literature Review}\label{sec-chapter2}

This chapter develops the theoretical and empirical foundation for the
proposed framework through a critical review of four interconnected
bodies of literature: the conceptual foundations of technical debt,
existing measurement and prioritisation approaches, technical debt
governance in financial and regulated enterprise IT, and multi-criteria
decision-making. Each section synthesises convergent findings,
identifies tensions and limitations, and positions these gaps in
relation to the research design. The review adopts a critical,
problem-oriented perspective rather than an exhaustive catalogue,
building a coherent argument that existing approaches are not adequate
for service-level enterprise IT contexts and that a quantitative
framework applying TOPSIS to operational service management data
represents a defensible and novel contribution.

\section{Context of the Topic}\label{context-of-the-topic}

Technical debt has become both a central concept and a persistent
management problem in software engineering and enterprise IT governance
(Avgeriou \emph{et al.}, 2023). Ward Cunningham first introduced the
metaphor in 1992, describing technical debt as the long-term cost of
short-term technical compromises (Cunningham, 1992). Over time, the
concept has expanded from a focus on poor code to a broader
understanding that also includes architectural, infrastructure,
organisational and process aspects of complex IT systems (Kruchten, Nord
and Ozkaya, 2012).

Despite this conceptual expansion, much of the empirical and
tool-supported literature remains anchored at the source code level,
relying heavily on static analysis tools and code-based metrics. These
approaches work tolerably well for greenfield development dominated by
in-house code but are much less appropriate for enterprise environments
where operational reliability depends on vendor platforms, legacy
infrastructure and tightly integrated service stacks (Amanatidis
\emph{et al.}, 2020).

The present review is motivated by a gap at the intersection of
technical debt management, multi-criteria decision-making and
mission-critical enterprise IT operations. To the best of current
literature, no framework provides a quantitative, service-level
mechanism for measuring and prioritising technical debt in environments
such as trading IT. In these contexts, significant forms of technical
debt manifest not only in code repositories but also in operational
outcomes such as service instability, patch delay and infrastructure
obsolescence (Ramasubbu and Kemerer, 2021; Haki \emph{et al.}, 2023).
The following sections build a theoretical and methodological basis for
the proposed framework.

\section{Methodology for Literature
Selection}\label{methodology-for-literature-selection}

The literature search was conducted between January and April 2026 using
a systematic and reproducible approach, drawing on four principal
academic databases: the University of Essex library, Scopus, IEEE Xplore
and Google Scholar. Search terms included the following.

\begin{longtable}[]{@{}
  >{\raggedright\arraybackslash}p{(\linewidth - 2\tabcolsep) * \real{0.3000}}
  >{\raggedright\arraybackslash}p{(\linewidth - 2\tabcolsep) * \real{0.7000}}@{}}
\caption{Search terms used for literature
review}\label{tbl-search-terms}\tabularnewline
\toprule\noalign{}
\begin{minipage}[b]{\linewidth}\raggedright
Theme
\end{minipage} & \begin{minipage}[b]{\linewidth}\raggedright
Search String
\end{minipage} \\
\midrule\noalign{}
\endfirsthead
\toprule\noalign{}
\begin{minipage}[b]{\linewidth}\raggedright
Theme
\end{minipage} & \begin{minipage}[b]{\linewidth}\raggedright
Search String
\end{minipage} \\
\midrule\noalign{}
\endhead
\bottomrule\noalign{}
\endlastfoot
TD foundations &
\texttt{"technical\ debt"\ AND\ (metaphor\ OR\ taxonomy\ OR\ typology\ OR\ theory)} \\
TD measurement and tools &
\texttt{"technical\ debt"\ AND\ (measurement\ OR\ "static\ analysis"\ OR\ SonarQube\ OR\ metric\ OR\ "code\ smell")} \\
TD prioritisation &
\texttt{"technical\ debt"\ AND\ (priorit*\ OR\ "backlog\ management"\ OR\ remediation\ OR\ repayment)} \\
TD in enterprise and financial IT &
\texttt{"technical\ debt"\ AND\ ("enterprise"\ OR\ "financial\ services"\ OR\ "trading"\ OR\ "legacy\ system"\ OR\ "regulated")} \\
TD automation &
\texttt{"technical\ debt\ management"\ AND\ ("machine\ learning"\ OR\ intelligent\ OR\ automation\ OR\ NLP)} \\
MCDM and TOPSIS &
\texttt{TOPSIS\ AND\ ("multi-criteria"\ OR\ MCDM\ OR\ "SAW"\ OR\ "decision\ making"\ OR\ priorit*)} \\
Patch management &
\texttt{"patch\ management"\ AND\ (priorit*\ OR\ risk\ OR\ vulnerability\ OR\ SLA\ OR\ enterprise)} \\
Operational risk and DORA &
\texttt{("DORA\ metrics"\ OR\ "mean\ time\ to\ restore")\ AND\ ("technical\ debt"\ OR\ "software\ delivery")} \\
\end{longtable}

Sources were screened against the following pre-specified inclusion and
exclusion criteria:

\begin{longtable}[]{@{}
  >{\raggedright\arraybackslash}p{(\linewidth - 2\tabcolsep) * \real{0.3000}}
  >{\raggedright\arraybackslash}p{(\linewidth - 2\tabcolsep) * \real{0.7000}}@{}}
\caption{Inclusion and exclusion
criteria}\label{tbl-criteria}\tabularnewline
\toprule\noalign{}
\begin{minipage}[b]{\linewidth}\raggedright
Criterion
\end{minipage} & \begin{minipage}[b]{\linewidth}\raggedright
Rule
\end{minipage} \\
\midrule\noalign{}
\endfirsthead
\toprule\noalign{}
\begin{minipage}[b]{\linewidth}\raggedright
Criterion
\end{minipage} & \begin{minipage}[b]{\linewidth}\raggedright
Rule
\end{minipage} \\
\midrule\noalign{}
\endhead
\bottomrule\noalign{}
\endlastfoot
Publication type & Peer-reviewed journal articles, conference
proceedings or empirical technical reports \\
Language & English full text available \\
Date range & 2019--2026, or foundational regardless of date \\
Thematic relevance & Directly addresses at least one of the four
thematic pillars \\
Exclusion: scope & Opinion pieces without empirical grounding \\
Exclusion: topic & Agile process improvement without a measurable TD
component \\
Exclusion: duplicates & Superseded by a more comprehensive journal
version \\
\end{longtable}

Screening was conducted in two stages. Titles and abstracts of all
retrieved records were first assessed against these criteria, reducing
the initial pool. Remaining sources then underwent full-text review to
verify relevance and methodological quality. Forward and backward
citation snowballing was subsequently applied to three anchor
publications: Kruchten, Nord and Ozkaya (2012), Lenarduzzi \emph{et al.}
(2021) and Forsgren, Humble and Kim (2018) to surface relevant studies
not captured by database searches. The process yielded approximately 110
candidate sources, of which more than 50 were retained following
full-text screening.

The review uses narrative synthesis rather than a formal systematic
review protocol. The research spans three disciplines, namely software
engineering, enterprise IT governance and decision science, which differ
significantly in their research methods, terminology and evidence
standards. A single meta-analytic framework would not accommodate this
heterogeneity, and the aim of the review is to build a coherent argument
in support of a novel design science artefact, for which narrative
synthesis is the appropriate approach (Hevner \emph{et al.}, 2004). The
following four sections build the theoretical and contextual case for
the proposed framework in sequence.

\section{Conceptual Foundations of Technical
Debt}\label{conceptual-foundations-of-technical-debt}

The first thematic area traces the evolution of the technical debt
concept from the original metaphor to later theoretical developments,
and then to its expansion beyond software into systems engineering,
infrastructure and vendor platforms. Understanding this evolution is
essential because most existing tools and practices are still tuned to
an earlier, narrower view of technical debt that misaligns with the
realities of contemporary enterprise IT.

\subsection{Origins and the Metaphor}\label{origins-and-the-metaphor}

Cunningham (1992) introduced the metaphor to explain his refactoring
strategy to non-technical stakeholders. In his original formulation,
shipping code that is ``not quite right'' resembles taking on financial
debt, with each subsequent maintenance effort representing interest. If
the principal is not repaid through refactoring, the cost of working
with the system compounds until change becomes very difficult.

The debt metaphor has proved productive, but it carries important
conceptual limits. In financial debt, the principal is contractually
defined, and the interest rate is known in advance, whereas technical
debt has neither. The principal is an estimated remediation cost that
changes as systems evolve. The interest depends on future development
activity that is uncertain at the time of incurrence. These limitations
matter directly for the present research because a central objective of
the proposed framework is to make interest cost concrete and
operationally measurable. The framework therefore does not rely on the
metaphor for rigour; it grounds its estimates in service-level
operational data rather than in abstract financial analogies.

Several influential practitioners refined the concept. McConnell (2008)
introduced a cost structure distinguishing principal, recurring interest
and compounding interest, and differentiated unintentional debt from
intentional debt incurred as a conscious trade-off. Fowler (2009)
proposed a quadrant distinguishing reckless from prudent debt and
deliberate from inadvertent debt. It should be noted that the quadrant
originates in a practitioner blog post rather than a peer-reviewed
publication; it is a useful heuristic for stakeholder communication but
should not be treated as a theoretically validated framework. Kruchten,
Nord and Ozkaya (2012) subsequently formalised a taxonomy covering code,
design, architectural, test, build and documentation debt on the basis
of a structured research programme, and this taxonomy underpins much of
the subsequent empirical work.

\subsection{Theoretical Development}\label{theoretical-development}

Empirical studies have confirmed and extended this view. Besker (2020)
reported that developers in industrial settings spend roughly 23\% of
their time dealing with the consequences of existing technical debt,
with architectural debt exerting the strongest negative influence on
productivity. Bedi and Kaur (2022) used regression analysis on
open-source repositories and found that commit frequency, lines of code,
code smells, and test coverage are significant predictors of technical
debt, whereas developer reputation metrics are not. This suggests that
technical debt is driven primarily by development practices rather than
by the perceived quality of individual contributors.

Ahmad \emph{et al.} (2026) approached the topic qualitatively,
interviewing experts from four international companies. They identified
five non-technical dimensions of debt accumulation in large-scale agile
environments, namely social dynamics, process, people, documentation and
requirements, with lack of communication emerging as the dominant
driver. This reinforces the view of technical debt as a sociotechnical
phenomenon that cannot be governed by technical measurement alone.

Vidoni, Codabux and Fard (2022) advanced a more radical argument through
a game-theoretic lens, proposing that in complex systems technical debt
cannot be completely removed and that its effects cannot be fully
controlled even when individual items are repaid. The concept of
``infinite technical debt'' challenges the assumption, implicit in many
repayment-focused frameworks, that debt is always recoverable given
sufficient effort. This argument has not yet been empirically validated
in enterprise IT settings and should be understood as a theoretical
provocation rather than an established finding. Its practical
implication for the present framework is that governance outputs should
surface replacement decisions where service-level data suggests debt has
crossed a threshold beyond which incremental repayment is unlikely to be
effective.

\subsection{Expanding the Conceptual
Perimeter}\label{expanding-the-conceptual-perimeter}

The conceptual perimeter of technical debt has progressively widened to
include broader systems and lifecycle concerns. Kleinwaks, Batchelor and
Bradley (2023) conducted a systematic review of technical debt in
systems engineering and found that requirements, architecture,
integration artefacts and verification artefacts can all accumulate
debt, and that many software-oriented approaches are poorly adapted to
multidisciplinary systems engineering lifecycles. Hanssen, Brataas and
Martini (2019) identified scalability debt as a distinct subtype
triggered by regulatory change: in an open banking case study, they
showed that architectural decisions, adequate before the EU Payment
Services Directive 2 (PSD2), became sources of debt once new scalability
requirements were imposed without architectural review. This illustrates
how external regulatory change can convert previously adequate designs
into liabilities, a pattern directly relevant to financial trading IT
where regulatory obligations are similarly dynamic. Yang \emph{et al.}
(2019) proposed a taxonomy of Commercial Off-the-Shelf (COTS) related
technical debt covering seven key types, including vendor platform
obsolescence and integration constraints. Managing such forms of debt
requires different tools from those used for internally developed code.
For financial trading IT, where vendor platforms and COTS products are
central, this broader view is particularly important: many of the
significant forms of technical debts, including platform debt,
infrastructure debt and vulnerability debt, are not detectable by
standard static analysis.

\section{Measurement and Prioritisation
Frameworks}\label{measurement-and-prioritisation-frameworks}

The second thematic area examines how technical debt is measured and
prioritised in practice. It argues that current measurement is heavily
concentrated at the code level and that existing prioritisation
frameworks lack the multidimensional quantitative rigour required for
robust governance in enterprise settings.

\subsection{The Dominance of Static Analysis and Its
Limits}\label{the-dominance-of-static-analysis-and-its-limits}

Static code analysis tools such as SonarQube, SQALE and CAST dominate
current practice, typically operationalising technical debt as a
weighted sum of estimated remediation effort derived from rule
violations (Biazotto \emph{et al.}, 2024). In a systematic mapping
study, Biazotto \emph{et al.} (2024) identified 178 primary studies and
found that most automation artefacts operate at the code level, that
substantial human intervention is still required even at higher
automation levels, and that end-to-end automation across the technical
debt management lifecycle is rare. The focus on code leaves
infrastructure-level and platform-level debt largely unaddressed.
Ampatzoglou \emph{et al.} (2022) presented SDK4ED as a comprehensive
industry oriented platform integrating several code level techniques. A
three-year evaluation suggested that SDK4ED can deliver financial
benefits to software organisations. At the same time, the platform
offers only limited support for debt related to infrastructure or vendor
platforms. Tornhill and Borg (2022) analysed 39 proprietary production
codebases and found that poor code quality is statistically associated
with higher defect density, longer remediation cycles and reduced
developer throughput, confirming that code-level debt has clear business
consequences. This does not change the fact that an exclusive focus on
code creates blind spots: organisations may allocate significant
resources to refactoring while leaving critical legacy infrastructure
and unsupported components unaddressed.

Amanatidis \emph{et al.} (2020) evaluated agreement between different
technical debt assessment tools applied to 50 open-source projects and
observed substantial divergence: tools applied to the same codebase
produce different estimates and highlight different violation sets, with
no single tool consistently dominating. This implies that a debt figure
generated by any one tool is a function of that tool's ruleset and
thresholds rather than an objective measurement. Albuquerque \emph{et
al.} (2023) reinforced this finding at scale, noting that infrastructure
debt and process debt remain significantly understudied compared to code
and architectural debt.

\subsection{The Prioritisation Gap}\label{the-prioritisation-gap}

Lenarduzzi \emph{et al.} (2021) carried out the most comprehensive
review of technical debt prioritisation strategies to date, analysing 44
primary studies. They found that existing approaches cluster around
cost-based, value-based and risk-based strategies, that most adopt a
single primary criterion, that quantitative methods are less common than
heuristic techniques, and that empirical validation on industrial
datasets is rare. They concluded that the field lacks agreement on which
factors matter, how to measure them and what constitutes an adequate
prioritisation process. This conclusion directly motivates the
quantitative contribution of the present work. The stakes of inadequate
prioritisation are concrete: de Toledo \emph{et al.} (2021) showed in a
production microservices system that targeted repayment of architectural
technical debt led to an 84\% reduction in incidents, anchoring the case
for prioritisation in a measurable operational outcome.

Empirical studies consistently reveal a gap between tool-generated
priorities and practitioner behaviour. Practitioners consider both value
and cost rather than a single strategy, and perceived value is highly
context-dependent (Alfayez \emph{et al.}, 2023). Pina, Seaman and
Goldman (2022) surveyed developers and found that prioritisation
decisions reflect a combination of technical severity, business impact
and available effort, and that developers frequently experience tension
between technically optimal repayment choices and organisationally
imposed constraints. Most IT managers nonetheless lack a formal
technical debt management process, with cross-functional communication
identified as the primary obstacle (Wiese and Borowa, 2023). Stochel
\emph{et al.} (2023) proposed the Continuous Debt Valuation Approach
(CoDVA), which assigns a financial value to each debt item based on
remediation cost, interest rate and business risk exposure; case studies
show that financial valuation can significantly change repayment
priorities relative to effort-only approaches, but CoDVA still relies on
development-level data and does not incorporate operational service
management signals such as incident rates or change failure rates.

These studies collectively indicate that any quantitative prioritisation
framework embedding value judgements about which systems and risks
matter most, must make those judgements transparent and auditable.
Particularly in regulated institutions, where priority scores may
influence funding, staffing and post-incident accountability.

\section{Technical Debt in Financial and Regulated Enterprise
IT}\label{technical-debt-in-financial-and-regulated-enterprise-it}

The third thematic area applies these general insights to the specific
setting of financial services and regulated trading IT, focusing on how
debt manifests operationally, how it translates into measurable harm,
and how related work in security and patch management quantifies that
harm.

\subsection{The Financial Services
Context}\label{the-financial-services-context}

Financial services organisations are characterised by complex, layered
legacies, strict regulatory compliance requirements and stringent
service availability obligations, with architectures that often depend
on vendor platforms and COTS products that cannot be freely refactored
(Haki \emph{et al.}, 2023).

Haki \emph{et al.} (2023) conducted clinical research at Credit Suisse,
designing and introducing a digital nudge for technical debt management
across more than 3,000 IT applications. As part of this work, they
defined a technical debt taxonomy with four main categories: code
quality debt, infrastructure lifecycle debt, vulnerability debt and
automation debt. Infrastructure lifecycle debt was measured as
cumulative months of unsupported components; vulnerability debt was
derived from weighted counts of scanner-detected vulnerabilities. This
taxonomy is significant for two reasons. First, it shows that a major
financial institution assigns equal operational importance to
infrastructure and vulnerability debt as to code-level debt. Second,
this directly supports the broader measurement approach taken in the
present research.

The empirical basis of the evaluation is, however, limited. The
quantitative analysis focused on three development teams, two of which
reduced their debt index over one year while the third did not. The
authors identified management intent as the decisive factor, rather than
the nudge design itself. The clinical research design, with researchers
embedded in the organisation, also raises questions about independence
and replicability. These factors do not invalidate the conceptual
contribution, but they caution against generalising from the observed
outcomes to other institutions or debt types.

Rolland and Lyytinen (2021) investigated the organisational dynamics of
architectural debt accumulation at a large Scandinavian financial
institution through eleven in-depth interviews with business managers,
enterprise architects and developers. Their analysis identified two
persistent tensions: between decentralised feature delivery teams
optimising for local throughput and the governance required to manage
debt at the system level, and between the pace of digital innovation and
the long-term sustainability of the supporting architecture. Both
tensions are intensified in financial services by compliance obligations
and competitive pressure to release new products rapidly. The study is
notable for grounding these dynamics in empirical data from a regulated
banking institution. Its primary limitation is the single-case design;
whether the patterns observed generalise to institutions with different
architectural profiles, including trading-specific environments, remains
an open question.

\subsection{Enterprise Systems and the Outsourcing
Context}\label{enterprise-systems-and-the-outsourcing-context}

Banker, Liang and Ramasubbu (2020) examined 26 firms operating a
commercial COTS CRM system over its full eleven-year lifecycle. Using
propensity score matching against firms with zero technical debt, they
found that a 10\% increase in technical debt reduced gross return on
assets by 16\% on average. Critically, this penalty more than doubled
within six years of adoption as code decay compounded structural
complexity. While all firms initially benefited from system
customisation, those carrying median debt levels saw the net value of
their system turn negative by year five, a pattern robust to alternative
performance measures including asset turnover, gross margin and
operating income.

In a later study, Ramasubbu and Kemerer (2021) extended this analysis to
1,824 outsourced remediation projects involving 424 client
organisations. They found that ``control balancing'' between client and
vendor improved remediation outcomes only when both parties had agreed a
long-term migration roadmap. In the absence of such a roadmap, increased
coordination actually harmed performance. For prioritisation frameworks,
this implies that ranking debt items without embedding them in a
strategic roadmap can be counterproductive.

Rinta-Kahila \emph{et al.} (2023) examined a global manufacturing firm's
COTS implementation and showed how staged legacy replacement, combined
with recurrent customisations to restore local flexibility, led to
self-reinforcing architectural debt accumulation. A system dynamics
model captured reinforcing feedback loops between customisation
decisions, debt accumulation and declining digital options. This work
suggests that organisations do not typically neglect technical debt
deliberately. Instead, they are trapped by interactions between
technical design choices and organisational routines that standard
code-level tools cannot reveal. Governance-level visibility of
infrastructure and architectural debt is required to detect these
dynamics before they become irreversible.

\subsection{DORA Metrics as Operational Debt
Proxies}\label{dora-metrics-as-operational-debt-proxies}

Forsgren, Humble and Kim (2018) validated DORA metrics across more than
2,000 organisations, showing that MTTR, CFR, deployment frequency and
lead time for changes distinguish high-performing from low-performing
delivery teams and can be derived from ITSM tools without source code
access. It is important to acknowledge that DORA metrics were validated
as delivery performance predictors, not as technical debt proxies. The
inferential step from degraded indicators to elevated TD severity is
plausible but not empirically established as a general relationship. The
present framework therefore treats this as a testable hypothesis rather
than a given premise. Nord, Ozkaya and Woody (2021) provide more direct
grounding by classifying delivery tempo degradation and deployment
instability explicitly as observable TD symptoms detectable from
production records, positioning MTTR and CFR as diagnostic signals of
accumulated debt rather than borrowed performance surrogates.

The evidence nonetheless requires careful qualification. Paudel \emph{et
al.} (2025) found that TD density explained between 5\% and 41\% of lead
time variance across six fintech components, with team structure and
component ownership acting as confounding factors. This reinforces the
case for multi-criteria measurement rather than reliance on any single
proxy. Nielsen and Skaarup (2022) offer complementary empirical support.
In a vendor-managed public sector portfolio without in-house development
resources, operational errors emerged shortly after TD-related resource
constraints were imposed on a shared service component, and unsupported
component versions were only detected through infrastructure migration
rather than routine monitoring. While the public sector context limits
direct transferability, the operational mechanism is consistent with the
vendor-dominated profile of the case organisation.

The following table maps each operational indicator to the technical
debt dimension it proxies and the literature that justifies its
inclusion:

\begin{longtable}[]{@{}
  >{\raggedright\arraybackslash}p{(\linewidth - 4\tabcolsep) * \real{0.2500}}
  >{\raggedright\arraybackslash}p{(\linewidth - 4\tabcolsep) * \real{0.3000}}
  >{\raggedright\arraybackslash}p{(\linewidth - 4\tabcolsep) * \real{0.4500}}@{}}
\caption{Mapping of operational indicators to technical debt dimensions
and literature justification}\label{tbl-indicators-lit}\tabularnewline
\toprule\noalign{}
\begin{minipage}[b]{\linewidth}\raggedright
Indicator
\end{minipage} & \begin{minipage}[b]{\linewidth}\raggedright
TD Dimension
\end{minipage} & \begin{minipage}[b]{\linewidth}\raggedright
Literature Justification
\end{minipage} \\
\midrule\noalign{}
\endfirsthead
\toprule\noalign{}
\begin{minipage}[b]{\linewidth}\raggedright
Indicator
\end{minipage} & \begin{minipage}[b]{\linewidth}\raggedright
TD Dimension
\end{minipage} & \begin{minipage}[b]{\linewidth}\raggedright
Literature Justification
\end{minipage} \\
\midrule\noalign{}
\endhead
\bottomrule\noalign{}
\endlastfoot
Incident frequency & Architectural debt & de Toledo \emph{et al.}
(2021); Ramasubbu and Kemerer (2021); Ramasubbu, Kemerer and Woodard
(2015); Nielsen and Skaarup (2022) \\
Mean Time to Restore (MTTR) & Infrastructure / platform debt & Forsgren,
Humble and Kim (2018); Nord, Ozkaya and Woody (2021); Paudel \emph{et
al.} (2025) \\
Change failure rate & Process debt & Forsgren, Humble and Kim (2018);
Nord, Ozkaya and Woody (2021); Paudel \emph{et al.} (2025) \\
Patch recency & Vulnerability debt & Tizio, Armellini and Massacci
(2023); Markkanen and Frantti (2023) \\
Unsupported component months & Infrastructure lifecycle debt & Haki
\emph{et al.} (2023) \\
\end{longtable}

Using operational service management data as a basis for technical debt
prioritisation raises important ethical questions, addressed in full in
Section~\ref{sec-ethics}.

\subsection{Security Debt and Patch
Management}\label{security-debt-and-patch-management}

Security-related technical debt is most visible in the literature on
patch management and vulnerability handling. Tizio, Armellini and
Massacci (2023) analysed 86 Advanced Persistent Threat campaigns dated
between 2008 and 2020, finding that enterprises delaying updates by one
month faced 4.9 times higher odds of compromise compared to immediate
patching, rising to 9.1 times at three months. Contrary to common
assumption, most campaigns exploited publicly known vulnerabilities
rather than zero-days, confirming that unpatched systems carry
measurable, quantifiable risk. These figures should be contextualised
carefully: the dataset spans all industries and predates regulatory
changes, including FINMA circular 2023/1 and the EU DORA regulation,
that have since accelerated patching cycles in financial services
specifically (\emph{{FINMA Circular} 2023/1 - {Operational} risks and
resilience -- banks}, 2024). The risk multipliers are therefore best
treated as indicative of the direction and order of magnitude of the
relationship rather than as directly applicable benchmarks for
contemporary trading IT.

Markkanen and Frantti (2023) reviewed current NIST guidance and the
academic literature on patch management planning. They argue that
uniform, time-based patching policies are increasingly inadequate given
the shrinking window between vulnerability disclosure and exploitation,
and that organisations should treat patching as a risk-based,
organisation-wide business priority rather than a purely technical
responsibility. Nurse (2025) framed patch prioritisation as a problem of
maintaining SLA compliance, noting that business partners and clients
drive patching decisions through supply chain agreements and contractual
requirements. This aligns patch management directly with the
service-level focus of the present research.

Kula \emph{et al.} (2019) studied release practices and dependency
management at ING, showing that keeping upstream dependencies current is
a persistent challenge in rapid-release banking environments and that
delays in dependency updates contribute directly to technical debt. This
finding is consistent with the broader argument, that technical debt
manifests in ways detectable through operational service management
data. The study is based on a single-organisation case study, however,
and may reflect ING-specific governance practices; whether the same
dependency dynamics apply in trading IT environments with different
release governance structures is not established.

Together, these studies illustrate that in financial trading IT,
significant forms of technical debt extend well beyond source code.
Regulatory change can convert previously adequate architectures into
sources of debt overnight (Hanssen, Brataas and Martini, 2019),
unsupported components accumulate measurable infrastructure lifecycle
debt (Haki \emph{et al.}, 2023), and service instability signals
accumulated debt that compounds with system growth (Ramasubbu and
Kemerer, 2021). Each of these manifestations is detectable through ITSM
and CMDB records rather than static analysis, establishing the data
foundation for the proposed framework.

\section{Multi-Criteria Decision-Making and
TOPSIS}\label{multi-criteria-decision-making-and-topsis}

The fourth thematic area argues that multi-criteria decision-making
(MCDM) methods are well suited to the technical debt prioritisation
problem and that TOPSIS is more appropriate for this context than
simpler alternatives.

\subsection{The Case for Multi-Criteria
Methods}\label{the-case-for-multi-criteria-methods}

Technical debt prioritisation involves balancing several criteria
simultaneously: remediation cost, business value, security risk,
operational impact and regulatory compliance. Practitioners weigh these
dimensions together rather than applying a single factor (Alfayez
\emph{et al.}, 2023). Lenarduzzi \emph{et al.} (2021) concluded that the
absence of a validated, quantitative, multi-criteria framework is a
central obstacle to evidence-based technical debt governance.

A weighted scoring matrix is the simplest alternative, but it assumes
that criteria are independent and commensurable across scales, neither
of which holds reliably when combining continuous cost metrics with
ordinal or binary compliance indicators. A risk matrix is categorical
and does not support fine-grained ranking of a large service portfolio.
TOPSIS addresses both limitations by normalising heterogeneous criteria
before aggregation and ranking alternatives by their geometric proximity
to the ideal solution, making it well suited to contexts where criteria
differ in type, scale and direction of preference (Ciardiello and
Genovese, 2023).

\subsection{TOPSIS: Mechanics, Applications and Critical
Limitations}\label{topsis-mechanics-applications-and-critical-limitations}

Hwang and Yoon (1981) introduced TOPSIS, defining the ideal solution as
the hypothetical alternative scoring best on all criteria simultaneously
and ranking alternatives by Euclidean proximity to this ideal. Zanakis
\emph{et al.} (1998) showed that TOPSIS produces stable rankings close
to the Pareto-optimal ordering across a range of problem structures,
supporting its use in prioritisation tasks under uncertainty.

TOPSIS constructs a normalised weighted decision matrix, identifies
ideal best and worst values per criterion, and ranks alternatives by
their closeness coefficient, the ratio of the distance to the ideal
worst to the sum of both distances (Albarak \emph{et al.}, 2022). A
well-known limitation is rank reversal: adding or removing alternatives
can shift the rankings of existing options (García-Cascales and Lamata,
2012). Yang (2020) proposed a normalisation adjustment that mitigates
this by defining reference points independently of the current
alternative set, which is particularly relevant for enterprise
portfolios where services are regularly added or decommissioned.

Three alternative MCDM methods warrant comparison. VIKOR identifies
compromise solutions between conflicting criteria rather than producing
a global ranking, making it less appropriate when the output required is
a ranked remediation list (Shekhovtsov and Sałabun, 2020). AHP derives
weights through pairwise comparisons but does not natively produce a
full portfolio ranking, and the comparison process grows burdensome as
the number of criteria increases (Ishak and Wanli, 2020). A weighted
scoring matrix is simpler still but assumes criteria are measured on
comparable scales, which does not hold when mixing continuous
operational metrics with categorical compliance indicators. TOPSIS
integrates normalisation, weighting and ranking into a single procedure,
making it the most efficient choice for generating a ranked service
portfolio. Where expert-derived weights are needed, AHP pairwise
comparisons can serve as an upstream step before feeding the results
into TOPSIS, a combination demonstrated by Ishak and Wanli (2020) in a
regulated operational context.

Albarak \emph{et al.} (2022) provide the most directly analogous
precedent, combining TOPSIS with portfolio theory to prioritise
technical debt and showing that it surfaced high-priority items likely
to be deferred under effort-only approaches. TOPSIS has also been
applied in other risk-intensive, regulated contexts (Zytoon, 2020; Khan
and Purohit, 2022), providing methodological precedents for its use
here. None of these studies, however, operates on live service
management data in a large enterprise setting, which is a defining
characteristic of the present research.

\section{Transferability Critique and Ethical
Dimensions}\label{sec-transfer}

Although the preceding sections provide strong motivation for the
proposed framework, none of the existing applications reviewed operates
on live operational records in large enterprise settings. Most rely on
researcher-defined scenarios, synthetic datasets or controlled
development artefacts. The proposed research differs by applying TOPSIS
to operational records such as incident logs, ITSM change data and
vulnerability scan outputs, which are known to have quality issues.
Incident classifications can vary between teams. Recorded change failure
rates can be influenced by documentation practices as much as by
technical quality. Vulnerability counts depend on scanner coverage and
configuration. These data limitations must be explicitly acknowledged.
The sensitivity of rankings to input variation must be studied before
outputs can be interpreted with confidence (Greco \emph{et al.}, 2019).

Using operational service management data as a basis for technical debt
prioritisation also raises important ethical questions. Such records are
generated and owned by operations teams, and development teams may not
have consented to their use in performance evaluations. A framework that
ranks services based on operational indicators may influence budgets,
staffing and reputational assessments in regulated institutions. These
considerations are addressed in full in the methodology chapter.

\section{Synthesis: Converging Evidence and the Research
Gap}\label{synthesis-converging-evidence-and-the-research-gap}

The four thematic sections together define the conceptual,
methodological, contextual and technical landscape within which this
research is situated.

\subsection{Cross-Pillar Convergence}\label{cross-pillar-convergence}

The preceding sections converge on a common conclusion: the current
technical debt literature is conceptually rich but methodologically
fragmented, and it remains poorly aligned with the operational realities
of enterprise IT.

The first pillar concerns scope. Technical debt has expanded beyond
source code to include architecture, infrastructure, vendor platforms,
and operational dependencies (Yang \emph{et al.}, 2019; Kleinwaks,
Batchelor and Bradley, 2023), yet the dominant measurement approaches
still concentrate on code-level artefacts (Biazotto \emph{et al.},
2024). This leaves important forms of debt, especially those that emerge
in regulated and vendor-heavy environments, only partially observable
(Rinta-Kahila \emph{et al.}, 2023).

The second pillar concerns prioritisation. Existing studies show that
technical debt is not a single dimension but a multi-criteria problem
involving cost, value, risk, and operational impact (Lenarduzzi \emph{et
al.}, 2021; Alfayez \emph{et al.}, 2023). However, the literature does
not provide a stable quantitative method that combines these dimensions
into a transparent ranking mechanism for service-level decision-making
(Pina, Seaman and Goldman, 2022).

The third pillar concerns consequences. The empirical literature
consistently shows that unmanaged technical debt produces measurable
harm, including increased incidents, slower recovery, reduced
productivity, higher remediation cost, and greater vulnerability
exposure (Banker, Liang and Ramasubbu, 2020; de Toledo \emph{et al.},
2021; Tornhill and Borg, 2022). These effects are not abstract, they
appear in operational outcomes that can be observed, compared, and
prioritised (Ramasubbu and Kemerer, 2021).

The fourth pillar concerns measurement opportunity. Operational service
data, including incident frequency, recovery time, and change failure
rate, offers a practical way to observe debt related symptoms in
environments where source code is inaccessible or incomplete. Patch
recency can similarly be operationalised as an infrastructure lifecycle
debt proxy using unsupported component data (Haki \emph{et al.}, 2023),
and delayed patching has been shown to produce measurable operational
risk (Tizio, Armellini and Massacci, 2023). These indicators do not
claim to measure technical debt directly. Rather, they serve as proxies
for the operational consequences of debt accumulation and therefore
provide a defensible basis for prioritisation in enterprise IT.

Taken together, these pillars reveal the research gap addressed by this
dissertation. There is no validated service level framework that
combines operational indicators into a single ranked view of technical
debt priority in trading IT. The present study responds to that gap by
developing and evaluating a TOPSIS based framework grounded in
operational data and supported by sensitivity analysis (Hevner \emph{et
al.}, 2004).

\subsection{Cross-Cutting Limitations}\label{cross-cutting-limitations}

Several limitations qualify these conclusions. The empirical base for
financial IT and trading systems remains relatively thin. Studies at
Credit Suisse, at BankAlpha (a pseudonym) and at ING each provide
valuable insights but are restricted to a single organisation (Kula
\emph{et al.}, 2019; Rolland and Lyytinen, 2021; Haki \emph{et al.},
2023). This scarcity of empirical evidence from trading IT environments
specifically represents a gap that the present research is positioned to
address, contributing an operational dataset from a regulated financial
environment that is currently absent from the technical debt literature.

Existing applications of TOPSIS to technical debt rely on controlled
datasets rather than live operational records (Albarak \emph{et al.},
2022; Khan and Purohit, 2022). The comparative work of Shekhovtsov and
Sałabun (2020) shows that method choice matters less than weight
specification when weights are uncertain, which underscores the
importance of transparent and justified weight determination over method
selection per se.

Behavioural studies emphasise that management intent, incentives and
communication patterns often exert a stronger influence on remediation
outcomes than technical tools (Besker, 2020; Wiese and Borowa, 2023;
Haki \emph{et al.}, 2023). Technical frameworks can support but not
replace governance and leadership.

\subsection{The Research Gap}\label{the-research-gap}

To the best of current literature, no framework integrates service-level
operational metrics, multi-criteria prioritisation and the specific
constraints of financial trading IT into a single, auditable technical
debt ranking instrument. Code-centric tools are structurally
inapplicable to vendor-hosted and regulated platforms. Single-criterion
models conflict with both the empirical evidence on practitioner
behaviour and the multidimensional nature of operational risk in trading
environments. The proposed framework addresses this gap directly. Its
evaluation will follow the approach established in adjacent quantitative
work, assessing ranking stability across weighting schemes and
validating indicator scores against historical operational records
(Greco \emph{et al.}, 2019; Albarak \emph{et al.}, 2022).

\bookmarksetup{startatroot}

\chapter{Research Design and Methodology}\label{sec-chapter3}

This chapter explains the philosophical approach, research strategy, and
design decisions that govern the construction and evaluation of the
proposed framework. It proceeds in two parts: the first covers the
research design from philosophical positioning through to the evaluation
approach; the second addresses the ethical and professional
considerations arising from the use of operational data in a regulated
financial institution.

\section{Research Design}\label{research-design}

This study is grounded in a pragmatist research philosophy. Pragmatism
holds that the value of a knowledge claim is determined by whether the
resulting artefact works in the context for which it was designed (Kelly
and Cordeiro, 2020). This position suits the present study because the
objective is not to establish universal laws about technical debt, but
to construct a quantitative framework useful for prioritisation
decisions in a specific enterprise IT environment. Technical debt is
treated here as a construct whose meaning is shaped by measurement
choices rather than as a directly observable quantity (Amanatidis
\emph{et al.}, 2020). The framework therefore grounds its estimates in
operationally defined proxies and acknowledges that the resulting scores
are representations, not facts.

\subsection{Research Strategy}\label{research-strategy}

Design Science Research (DSR) governs this study, as formalised by
Hevner \emph{et al.} (2004). DSR is appropriate here because the primary
output is an artefact rather than a theory: a TOPSIS-based
prioritisation framework operating on live ITSM data. The framework is
designed to address the gap identified in Chapter~\ref{sec-chapter2},
where no existing approach integrates service-level operational metrics
for technical debt prioritisation in trading IT. Lenarduzzi \emph{et
al.} (2021) confirm the broader absence of validated, quantitative
prioritisation frameworks across the field; the service-level and
enterprise-IT dimension of that gap is established through the synthesis
in Chapter~\ref{sec-chapter2}.

Hevner \emph{et al.} (2004) specify seven guidelines for rigorous DSR,
mapped to this study in Table~\ref{tbl-dsr-guidelines}. On one side sits
the problem space: the absence of a service-level TD framework. On the
other sits the solution space: the TOPSIS artefact. These are connected
through input knowledge, such as DORA metrics and MCDM methods, and
through research processes, including indicator selection and
sensitivity analysis. This structure positions the framework as
prescriptive design knowledge applicable to regulated enterprise
contexts.

\begin{longtable}[]{@{}
  >{\raggedright\arraybackslash}p{(\linewidth - 2\tabcolsep) * \real{0.3000}}
  >{\raggedright\arraybackslash}p{(\linewidth - 2\tabcolsep) * \real{0.7000}}@{}}
\caption{DSR guidelines and their
implementation}\label{tbl-dsr-guidelines}\tabularnewline
\toprule\noalign{}
\begin{minipage}[b]{\linewidth}\raggedright
DSR Guideline
\end{minipage} & \begin{minipage}[b]{\linewidth}\raggedright
Implementation in this Study
\end{minipage} \\
\midrule\noalign{}
\endfirsthead
\toprule\noalign{}
\begin{minipage}[b]{\linewidth}\raggedright
DSR Guideline
\end{minipage} & \begin{minipage}[b]{\linewidth}\raggedright
Implementation in this Study
\end{minipage} \\
\midrule\noalign{}
\endhead
\bottomrule\noalign{}
\endlastfoot
Design as an artefact & The TOPSIS-based prioritisation framework \\
Problem relevance & Established in Chapter~\ref{sec-chapter2}: no
existing framework integrates service-level operational metrics for TD
prioritisation in trading IT \\
Design evaluation & Sensitivity analysis; criterion validity against
historical operational records \\
Research contributions & Novel indicator set; first application of
TOPSIS to live ITSM data in financial trading IT \\
Research rigour & Systematic indicator selection; reproducible TOPSIS
computation; documented normalisation \\
Design as a search process & Iterative refinement of indicator selection
(literature candidates audited against ITSM/CMDB availability) and
weight assignment \\
Communication of research & This dissertation; practitioner-accessible
scoring output \\
\end{longtable}

The process follows Vom Brocke, Hevner and Maedche (2020)'s six-step
model: problem identification and motivation (Chapter~\ref{sec-chapter1}
and Chapter~\ref{sec-chapter2}), definition of objectives
(Chapter~\ref{sec-chapter2}), design and development of the artefact
(Chapter~\ref{sec-chapter4}), demonstration, evaluation
(Chapter~\ref{sec-chapter5}), and communication
(Chapter~\ref{sec-chapter6}). The sequence is not strictly linear;
indicator selection and evaluation criteria were refined iteratively as
constraints in the available data emerged during development.

\begin{figure}

\centering{

\includegraphics[width=6in,height=2.67in]{index_files/figure-latex/mermaid-figure-1.png}

}

\caption{\label{fig-dsr-grid}DSR Grid for the present study, adapted
from vom Brocke, Hevner and Maedche (2020)}

\end{figure}%

\subsection{Research Design}\label{research-design-1}

The study is structured as a single-organisation embedded case study,
situated within the trading IT division of one of Switzerland's three
largest universal banks. Turnbull, Chugh and Luck (2021) define this
design as an investigation into a single instance of observable complex
phenomena, constrained by clearly identifiable boundaries. Discrete
units of analysis, in this case individual IT services, are examined
within that bounded organisational setting.

Three factors make this design appropriate. First, the regulatory,
operational, and vendor-platform characteristics specific to Swiss
financial IT resist generalisation from software engineering studies
conducted elsewhere. Second, access to granular ITSM records is
unavailable through surveys or experiments. Third, Hevner \emph{et al.}
(2004) specify that a design artefact must be evaluated through
rigorous, well-executed methods. The present study satisfies this
requirement by evaluating the framework against operational data from
the environment for which it was designed.

The study adopts the individual IT service as the unit of analysis,
defined as a distinct configuration item in the organisation's ITSM
instance. In this vendor-heavy enterprise environment, source-code
access is unavailable for both outsourced components and on-premises
vendor software (Ramasubbu and Kemerer, 2021). Service-level operational
indicators are therefore tested as candidate proxies for technical debt
prioritisation, rather than assumed as established substitutes. The
study is cross-sectional, analysing a fixed observation window rather
than tracking changes over time, consistent with the objective of
demonstrating ranking validity at a point in time. Longitudinal
evaluation is identified as a direction for future work.

\subsection{Data Sources and Collection}\label{sec-datasources}

The literature review (Chapter~\ref{sec-chapter2}) identified candidate
indicators mapping technical debt dimensions to operational data (see
Table~\ref{tbl-indicators-lit}). A data audit against available ITSM and
CMDB fields refined this set prior to data collection. Documentation
completeness was excluded as no structured ITSM or CMDB field captures
it consistently across services; relevant records reside in unstructured
repositories not amenable to uniform extraction. Test coverage ratio and
architectural coupling were excluded because both require access to
assets unavailable in this environment: test coverage is derived from
CI/CD build artefacts or source repositories, and architectural coupling
is measured through static dependency graph analysis of compiled
code(Amanatidis \emph{et al.}, 2020); neither is accessible for
vendor-managed nor outsourced components. Raw vulnerability density was
excluded as subsumed by patch recency, which provides a more stable
CMDB-derived proxy grounded in regulatory patching compliance rather
than raw scanner counts whose coverage varies across components.

The five retained indicators are fully operationalisable from live
records held in the organisation's ITSM instance and CMDB. Records were
accessed under an agreed arrangement with the Head of Trading IT over a
twelve-month observation window (January--December 2025) and anonymised
at source, with service names replaced by alphanumeric identifiers and
all team and vendor labels removed before analysis.

\begin{longtable}[]{@{}
  >{\raggedright\arraybackslash}p{(\linewidth - 8\tabcolsep) * \real{0.2000}}
  >{\raggedright\arraybackslash}p{(\linewidth - 8\tabcolsep) * \real{0.4000}}
  >{\raggedright\arraybackslash}p{(\linewidth - 8\tabcolsep) * \real{0.1000}}
  >{\raggedright\arraybackslash}p{(\linewidth - 8\tabcolsep) * \real{0.1500}}
  >{\raggedright\arraybackslash}p{(\linewidth - 8\tabcolsep) * \real{0.1500}}@{}}
\caption{Operational definition of indicators used in the TOPSIS
framework}\label{tbl-indicators-def}\tabularnewline
\toprule\noalign{}
\begin{minipage}[b]{\linewidth}\raggedright
Indicator
\end{minipage} & \begin{minipage}[b]{\linewidth}\raggedright
Operational Definition
\end{minipage} & \begin{minipage}[b]{\linewidth}\raggedright
Source
\end{minipage} & \begin{minipage}[b]{\linewidth}\raggedright
Unit
\end{minipage} & \begin{minipage}[b]{\linewidth}\raggedright
Direction
\end{minipage} \\
\midrule\noalign{}
\endfirsthead
\toprule\noalign{}
\begin{minipage}[b]{\linewidth}\raggedright
Indicator
\end{minipage} & \begin{minipage}[b]{\linewidth}\raggedright
Operational Definition
\end{minipage} & \begin{minipage}[b]{\linewidth}\raggedright
Source
\end{minipage} & \begin{minipage}[b]{\linewidth}\raggedright
Unit
\end{minipage} & \begin{minipage}[b]{\linewidth}\raggedright
Direction
\end{minipage} \\
\midrule\noalign{}
\endhead
\bottomrule\noalign{}
\endlastfoot
Incident frequency & Count of P2 and P3 incidents per service within the
observation window, normalised by service age in months. P1 incidents,
defined by the organisation as cases where the majority of its business
is affected, did not occur within the observation period and are
excluded. & ITSM & Incidents / month & Cost \\
MTTR & Mean elapsed time from incident creation to resolution status,
averaged across all P2 and P3 incidents per service & ITSM & Hours &
Cost \\
CFR & Proportion of change records per service that resulted in an
unplanned incident or rollback within 72 hours of deployment & ITSM &
Percentage (0--100) & Cost \\
Patch recency & Proportion of components within each service's
dependency stack recorded as within a vendor-supported version window at
point of extraction & CMDB & Percentage (0--100) & Benefit \\
Unsupported component months & Cumulative months of exposure to & & & \\
unsupported components within the service's CMDB dependency stack during
the observation window & CMDB & Months & Cost & \\
\end{longtable}

Incident frequency, MTTR and CFR proxy the operational consequences of
architectural and infrastructure debt. Service instability and slow
recovery are empirically linked to TD severity in vendor-managed
portfolios (Ramasubbu and Kemerer, 2021; Nielsen and Skaarup, 2022).
Deployment instability is classified as a direct observable symptom of
accumulated debt (Forsgren, Humble and Kim, 2018; Nord, Ozkaya and
Woody, 2021). Nord, Ozkaya and Woody (2021) further ground the inclusion
of MTTR and CFR specifically by classifying delivery tempo degradation
as a diagnostic TD signal detectable from production records,
independently of source code access. Patch recency captures the
proportion of components currently within a vendor-supported version
window. Unsupported component months measures the cumulative duration of
exposure to end-of-life components during the observation period.
Together, these two CMDB-derived indicators operationalise
infrastructure lifecycle and vulnerability debt following Haki \emph{et
al.} (2023), differentiating services with similar current patch status
but differing historical debt accumulation. Service-level operational
indicators are treated as candidate proxies for technical debt
prioritisation rather than established substitutes. The validity of this
inference is the central empirical hypothesis tested in
Chapter~\ref{sec-chapter5}.

\subsection{Framework Design}\label{sec-frameworkdesign}

The analytical core is the Technique for Order of Preference by
Similarity to Ideal Solution (TOPSIS), formulated by Hwang and Yoon
(1981) and validated through comprehensive reviews of its enterprise
applications (Pandey, Komal and Dincer, 2023). TOPSIS addresses the
single-criterion prioritisation limitation identified in
Chapter~\ref{sec-chapter2} by ranking services according to their
geometric proximity to ideal solutions after normalisation and
weighting. Given a decision matrix with \(m\) services and \(n\)
indicators, the method proceeds as follows.

The matrix is normalised using vector normalisation, which scales each
value relative to the column norm to ensure comparability across
indicators with different units and ranges (Pandey, Komal and Dincer,
2023):

\[
r_{ij} = \frac{x_{ij}}{\sqrt{\sum_{i=1}^{m} x_{ij}^2}}
\]

A weighted normalised matrix is constructed using criterion weight
\(w_j\):

\[
v_{ij} = w_j \cdot r_{ij}
\]

The positive and negative ideal solutions are identified:

\[
A^{\text{+}} = \{ v_1^{\text{+}}, \ldots, v_n^{\text{+}} \}, \quad A^{\text{-}} = \{ v_1^{\text{-}}, \ldots, v_n^{\text{-}} \}
\]

Euclidean distances from each ideal are calculated:

\[
S_i^{\text{+}} = \sqrt{\sum_{j=1}^{n}(v_{ij} - v_j^{\text{+}})^2}, \quad
S_i^{\text{-}} = \sqrt{\sum_{j=1}^{n}(v_{ij} - v_j^{\text{-}})^2}
\]

The relative closeness coefficient \(C_i\) constitutes the final score:

\[
C_i = \frac{S_i^{\text{-}}}{S_i^{\text{+}} + S_i^{\text{-}}}, \quad C_i \in [0, 1]
\]

A higher \(C_i\) indicates a service closer to the worst-case
operational profile, representing a higher technical debt priority.
Incident frequency, MTTR, and CFR are cost criteria where lower values
are preferable; patch recency is a benefit criterion where higher values
are preferable.

Figure~\ref{fig-framework-pipeline} illustrates the full pipeline from
indicator collection through to ranked output.

\begin{figure}

\centering{

\includegraphics[width=6.5in,height=0.95in]{index_files/figure-latex/mermaid-figure-3.png}

}

\caption{\label{fig-framework-pipeline}Framework pipeline: from
operational indicators to ranked prioritisation}

\end{figure}%

Two weight configurations are used. A baseline equal-weight
configuration (\(w_j = 0.25\)) serves as the reference case, in
accordance with Greco \emph{et al.} (2019). A secondary configuration
applies differential weights reflecting the higher operational salience
of MTTR and CFR in a trading environment, as established in
Chapter~\ref{sec-chapter2}. The sensitivity of rankings to weight
perturbation is assessed in Chapter~\ref{sec-chapter5}.

\subsection{Evaluation Approach}\label{evaluation-approach}

The framework is evaluated against a single criterion following the
sensitivity practice recommended by Greco \emph{et al.} (2019) for
composite indicators. Ranking stability is assessed by perturbing each
criterion weight by plus or minus 20 percentage points, with remaining
weights adjusted proportionally, and measuring rank-order correlations
between the baseline ranking and each perturbed variant using Kendall's
\(\tau\). This constitutes an analytical evaluation of the artefact in
the sense of Hevner \emph{et al.} (2004) Guideline 3, which requires
that the utility and quality of a design artefact be rigorously
demonstrated via well-executed evaluation methods. A consistently high
\(\tau\) across perturbations provides evidence that the ranking is not
sensitive to reasonable variation in weight assumptions, which is a
necessary condition for the framework to be actionable in practice.

\section{Ethical and Professional Considerations}\label{sec-ethics}

The ethical considerations arising from this study concern data
confidentiality and legal compliance, the responsible use of
prioritisation outputs, and alignment with professional standards.

\subsection{Data Confidentiality and
Anonymisation}\label{data-confidentiality-and-anonymisation}

The data originate from a live Swiss financial institution that is not
identified by name. All service, team, and vendor identifiers were
replaced with alphanumeric codes before analysis. No employee data are
present: all five indicators are service-level aggregates that cannot be
attributed to individuals. The data-sharing arrangement was established
with the explicit agreement of the Head of Trading IT, access was
limited to the required records, and no data were retained beyond the
analytical period or shared with third parties.

\subsection{Legal Compliance}\label{legal-compliance}

The organisation is subject to the revised Federal Act on Data
Protection (nFADP), which entered into force on 1 September 2023
(\emph{Federal {Act} on {Data Protection}}, 2023), enforced by the
Federal Data Protection and Information Commissioner (FDPIC). Because
the dataset contains no personal data, the nFADP's provisions on
personal data processing do not apply directly. The principle of data
minimisation, which the nFADP enshrines in alignment with international
standards, was nonetheless observed: only the metrics required for the
five TOPSIS indicators were extracted.

As a systemically significant institution, the organisation is also
subject to FINMA Circular 2023/1 on operational risks and resilience, in
force since 1 January 2024 (\emph{{FINMA Circular} 2023/1 -
{Operational} risks and resilience -- banks}, 2024). This circular
requires comprehensive ICT operational risk management, including
incident monitoring and service resilience. Areas structurally aligned
with the framework's operational focus. Although the specific metrics
used here are organisationally defined rather than prescribed by FINMA.
No ethics committee approval was required under the University of
Essex's research ethics framework, and a departmental no-objection
confirmation was obtained.

\subsection{Responsible Use of Prioritisation
Outputs}\label{responsible-use-of-prioritisation-outputs}

A relative rank assigned to a service carries a risk of being misused to
evaluate team performance. High-debt service profiles typically reflect
historical design decisions or legacy constraints rather than the
competence of the team currently managing the service (Ahmad \emph{et
al.}, 2026). Three mitigations address this risk: framework
documentation explicitly states that \(C_i\) reflects the operational
state of a service rather than team performance; outputs are intended
for IT governance and portfolio management functions only; and the
sensitivity analysis makes weighting assumptions transparent, reducing
the risk of false precision.

\subsection{Professional Standards and
Limitations}\label{professional-standards-and-limitations}

This study was conducted in accordance with the BCS Code of Conduct
(\emph{Code of {Conduct} for {BCS Members}}, 2022), including the
obligations to act with integrity, to respect the rights of third
parties, and to ensure that the artefact is fit for purpose. Three
limitations clarify scope. The single-site design means that indicator
weights and TOPSIS configuration are calibrated for one organisation and
should not be transferred without revalidation. The cross-sectional
window means the framework captures a snapshot; recently remediated
services may be under-ranked. Finally, the five indicators do not
exhaust all dimensions of operational technical debt; indicators related
to documentation, test coverage, or architectural coupling were excluded
due to data availability and remain candidates for future iterations.

\section{Rework!}\label{rework}

The following diagram illustrates the logical pipeline of the proposed
framework, from the collection of operational indicators through to a
ranked prioritisation output.

\begin{figure}

\centering{

\includegraphics[width=6in,height=0.42in]{index_files/figure-latex/mermaid-figure-2.png}

}

\caption{\label{fig-framework}Proposed framework pipeline}

\end{figure}%

\bookmarksetup{startatroot}

\chapter{Framework Development and Results}\label{sec-chapter4}

This chapter presents the construction and application of the
TOPSIS-based prioritisation framework. It covers data preparation,
indicator normalisation and weight assignment, TOPSIS computation and
ranking, and sensitivity analysis, before discussing the substantive
interpretation of the results. The complete Python implementation is
provided in Appendix A. All code referenced in this chapter can be
executed directly to reproduce the outputs shown here.

\section{Data Preparation}\label{data-preparation}

The dataset covers five IT services drawn from the trading IT division
of the case organisation, extracted from its ITSM instance. Services
were included on two criteria: they must be registered as active
configuration items and must have a minimum operational history of six
months within the twelve-month observation window, sufficient to produce
stable indicator estimates. Services with fewer than three recorded
production changes during the window were assigned the portfolio mean
for CFR, consistent with the imputation approach recommended for sparse
operational datasets (Greco \emph{et al.}, 2019). All service
identifiers were replaced with alphanumeric codes prior to analysis, as
described in Section 3.2.1. No team, vendor or individual identifiers
are present in the dataset.

The four indicators are operationally defined as follows, consistent
with the definitions established in Section 3.1.4. Incident frequency is
the total number of P2 and P3 incidents attributed to the service during
the observation window, divided by service age in months. MTTR is the
arithmetic mean of elapsed hours between incident creation and
resolution across all P2 and P3 incidents within the window. CFR is the
proportion of production changes that resulted in an unplanned rollback
or a subsequent incident within 72 hours of deployment. Patch currency
is the proportion of components in the service's CMDB dependency stack
whose installed version falls within a currently supported release
window.

\begin{longtable}[]{@{}
  >{\raggedright\arraybackslash}p{(\linewidth - 8\tabcolsep) * \real{0.2000}}
  >{\raggedright\arraybackslash}p{(\linewidth - 8\tabcolsep) * \real{0.2000}}
  >{\raggedright\arraybackslash}p{(\linewidth - 8\tabcolsep) * \real{0.2000}}
  >{\raggedright\arraybackslash}p{(\linewidth - 8\tabcolsep) * \real{0.2000}}
  >{\raggedright\arraybackslash}p{(\linewidth - 8\tabcolsep) * \real{0.2000}}@{}}
\caption{Raw decision matrix}\label{tbl-raw-data}\tabularnewline
\toprule\noalign{}
\begin{minipage}[b]{\linewidth}\raggedright
Service
\end{minipage} & \begin{minipage}[b]{\linewidth}\raggedright
Incident Freq.
\end{minipage} & \begin{minipage}[b]{\linewidth}\raggedright
MTTR (hrs)
\end{minipage} & \begin{minipage}[b]{\linewidth}\raggedright
CFR
\end{minipage} & \begin{minipage}[b]{\linewidth}\raggedright
Patch Recency
\end{minipage} \\
\midrule\noalign{}
\endfirsthead
\toprule\noalign{}
\begin{minipage}[b]{\linewidth}\raggedright
Service
\end{minipage} & \begin{minipage}[b]{\linewidth}\raggedright
Incident Freq.
\end{minipage} & \begin{minipage}[b]{\linewidth}\raggedright
MTTR (hrs)
\end{minipage} & \begin{minipage}[b]{\linewidth}\raggedright
CFR
\end{minipage} & \begin{minipage}[b]{\linewidth}\raggedright
Patch Recency
\end{minipage} \\
\midrule\noalign{}
\endhead
\bottomrule\noalign{}
\endlastfoot
SVC-01 & 1.8 & 38.4 & 0.27 & 0.52 \\
SVC-02 & 0.9 & 15.2 & 0.12 & 0.75 \\
SVC-03 & 2.4 & 47.1 & 0.33 & 0.46 \\
SVC-04 & 0.3 & 5.8 & 0.03 & 0.94 \\
SVC-05 & 0.4 & 44.2 & 0.05 & 0.89 \\
\end{longtable}

Incident frequency, MTTR and CFR are cost criteria, for which lower
values indicate lower technical debt severity. Patch currency is a
benefit criterion, for which higher values are preferable.

\subsection{Data Limitations}\label{data-limitations}

As established in Section~\ref{sec-transfer}, the operational data used
here carry inherent quality limitations that must be explicitly
acknowledged. Incident P2/P3 classifications vary across teams; CFR
reflects documentation practices as much as technical outcomes; and
patch recency depends on CMDB coverage and scanner configuration. These
do not invalidate the proxy approach but constrain score precision.

Mitigations include: (i) sensitivity analysis
(Section~\ref{sec-sensanaly}) to test ranking stability; and (ii)
cross-validation against independent remediation records
(Chapter~\ref{sec-chapter5}).

\section{Indicator Normalisation and
Weighting}\label{indicator-normalisation-and-weighting}

Before aggregation, each indicator is normalised using vector
normalisation to remove differences in scale and unit:

\[r_{ij} = \frac{x_{ij}}{\sqrt{\sum_{i=1}^{m} x_{ij}^2}}\]

Vector normalisation is preferred over min-max normalisation because it
preserves proportional relationships between values without compressing
the range to a fixed interval, a property that (Greco \emph{et al.},
2019) recommend for composite indicators where extreme values carry
substantive meaning. A service with an unusually high MTTR should retain
its full distance from lower-performing services in the normalised space
rather than be pulled toward a truncated maximum. A weighted normalised
matrix is then constructed by multiplying each normalised value by its
criterion weight:

\[v_{ij} = w_j \cdot r_{ij}\]

Two weight configurations are applied. The baseline equal-weight
configuration assigns \(w_j = 0.25\) to each criterion and serves as the
reference case against which the differential configuration is
evaluated. The differential configuration assigns higher weights to MTTR
and CFR.

This reflects a design judgement rather than a claim attributable to a
single empirical source, and it is grounded in three converging
observations from the literature. First, Ramasubbu, Kemerer and Woodard
(2015) establish that technical debt in COTS enterprise systems
increases the probability of system failures by up to 62 percent,
identifying incident frequency as the most directly observable TD signal
in vendor-platform environments. This makes the speed of recovery from
those incidents, captured by MTTR, the dimension most consequential to
operational governance once a failure has occurred. Second, Nord, Ozkaya
and Woody (2021) classify delivery tempo degradation and deployment
instability explicitly as observable symptoms from which technical debt
can be detected and prioritised in production environments (Nord, Ozkaya
and Woody, 2021). This provides direct grounding for elevating the
weights of both MTTR and CFR relative to the remaining two criteria.
Third, Fungprasertkul, Bahsoon and Kazman (2024) establish that TD
dimensions can and should be assigned differential weights to reflect
the perspectives of different stakeholders and the operational context
in which monitoring takes place (Fungprasertkul, Bahsoon and Kazman,
2024), offering methodological precedent for a two-configuration design
in which the baseline equal-weight case is tested against a
configuration calibrated to domain context.

Two weight configurations are applied, as specified in
Section~\ref{sec-frameworkdesign}. The baseline equal-weight
configuration assigns \(w_j = 0.25\) to each criterion and serves as the
reference case. The differential configuration assigns higher weights to
MTTR and CFR, reflecting their higher operational salience in a trading
environment where speed of recovery and deployment stability directly
determine financial exposure (Forsgren, Humble and Kim, 2018). The
differential weights are \(w_{\text{MTTR}} = w_{\text{CFR}} = 0.30\) and
\(w_{\text{IncFreq}} = w_{\text{Patch}} = 0.20\).

\begin{longtable}[]{@{}
  >{\raggedright\arraybackslash}p{(\linewidth - 8\tabcolsep) * \real{0.3200}}
  >{\raggedright\arraybackslash}p{(\linewidth - 8\tabcolsep) * \real{0.1700}}
  >{\raggedright\arraybackslash}p{(\linewidth - 8\tabcolsep) * \real{0.1700}}
  >{\raggedright\arraybackslash}p{(\linewidth - 8\tabcolsep) * \real{0.1700}}
  >{\raggedright\arraybackslash}p{(\linewidth - 8\tabcolsep) * \real{0.1700}}@{}}
\caption{Weight configurations}\label{tbl-weights}\tabularnewline
\toprule\noalign{}
\begin{minipage}[b]{\linewidth}\raggedright
Configuration
\end{minipage} & \begin{minipage}[b]{\linewidth}\raggedright
Incident Freq.
\end{minipage} & \begin{minipage}[b]{\linewidth}\raggedright
MTTR
\end{minipage} & \begin{minipage}[b]{\linewidth}\raggedright
CFR
\end{minipage} & \begin{minipage}[b]{\linewidth}\raggedright
Patch Currency
\end{minipage} \\
\midrule\noalign{}
\endfirsthead
\toprule\noalign{}
\begin{minipage}[b]{\linewidth}\raggedright
Configuration
\end{minipage} & \begin{minipage}[b]{\linewidth}\raggedright
Incident Freq.
\end{minipage} & \begin{minipage}[b]{\linewidth}\raggedright
MTTR
\end{minipage} & \begin{minipage}[b]{\linewidth}\raggedright
CFR
\end{minipage} & \begin{minipage}[b]{\linewidth}\raggedright
Patch Currency
\end{minipage} \\
\midrule\noalign{}
\endhead
\bottomrule\noalign{}
\endlastfoot
Equal (baseline) & 0.25 & 0.25 & 0.25 & 0.25 \\
Differential & 0.20 & 0.30 & 0.30 & 0.20 \\
\end{longtable}

\section{TOPSIS Computation and
Rankings}\label{topsis-computation-and-rankings}

The TOPSIS procedure follows the formulation of (Hwang and Yoon, 1981).
After constructing the weighted normalised matrix, the positive and
negative ideal solutions are:

\section{Sensitivity Analysis}\label{sec-sensanaly}

\section{Interpretation and
Discussion}\label{interpretation-and-discussion}

\bookmarksetup{startatroot}

\chapter{Evaluation}\label{sec-chapter5}

\section{Evaluation Aim}\label{evaluation-aim}

\section{Evaluation Criteria}\label{evaluation-criteria}

\section{Stability Results}\label{stability-results}

\section{Practical Validity}\label{practical-validity}

\section{Limitations of the
Evaluation}\label{limitations-of-the-evaluation}

\bookmarksetup{startatroot}

\chapter{Conclusion}\label{sec-chapter6}

\section{Research Summary}\label{research-summary}

\section{Findings Against Research
Questions}\label{findings-against-research-questions}

\section{Contribution to Knowledge}\label{contribution-to-knowledge}

\section{Practical Contribution}\label{practical-contribution}

\section{Limitations and Future Work}\label{limitations-and-future-work}

\section{Reflective Note}\label{reflective-note}

\bookmarksetup{startatroot}

\chapter*{References}\label{references}
\addcontentsline{toc}{chapter}{References}

\markboth{References}{References}

\phantomsection\label{refs}
\begin{CSLReferences}{0}{1}
\bibitem[\citeproctext]{ref-ahmadTechnicalDebtNot2026}
Ahmad, M.O., Mandić, V., Taušan, N., Katin, A. and Herath, P. (2026)
\textquotesingle Technical debt is not just technical: {An} industrial
case study in large agile software development\textquotesingle,
\emph{Journal of Systems and Software}, 234, p. 112719. Available at:
\url{https://doi.org/10.1016/j.jss.2025.112719}.

\bibitem[\citeproctext]{ref-albarakManagingTechnicalDebt2022}
Albarak, M., Bahsoon, R., Ozkaya, I. and Nord, R.L. (2022)
\textquotesingle Managing {Technical Debt} in {Database
Normalization}\textquotesingle, \emph{IEEE Transactions on Software
Engineering}, 48(3), pp. 755--772. Available at:
\url{https://doi.org/10.1109/TSE.2020.3001339}.

\bibitem[\citeproctext]{ref-albuquerqueManagingTechnicalDebt2023}
Albuquerque, D., Guimarães, E., Tonin, G., Rodríguezs, P., Perkusich,
M., Almeida, H., Perkusich, A. and Chagas, F. (2023)
\textquotesingle Managing {Technical Debt Using Intelligent Techniques}
- {A Systematic Mapping Study}\textquotesingle, \emph{IEEE Transactions
on Software Engineering}, 49(4), pp. 2202--2220. Available at:
\url{https://doi.org/10.1109/TSE.2022.3214764}.

\bibitem[\citeproctext]{ref-alfayezHowSonarQubeidentifiedTechnical2023}
Alfayez, R., Winn, R., Alwehaibi, W., Venson, E. and Boehm, B. (2023)
\textquotesingle How {SonarQube-identified} technical debt is
prioritized: {An} exploratory case study\textquotesingle,
\emph{Information and Software Technology}, 156, p. 107147. Available
at: \url{https://doi.org/10.1016/j.infsof.2023.107147}.

\bibitem[\citeproctext]{ref-amanatidisEvaluatingAgreementTechnical2020}
Amanatidis, T., Mittas, N., Moschou, A., Chatzigeorgiou, A.,
Ampatzoglou, A. and Angelis, L. (2020) \textquotesingle Evaluating the
agreement among technical debt measurement tools: Building an empirical
benchmark of technical debt liabilities\textquotesingle, \emph{Empirical
Software Engineering}, 25(5), pp. 4161--4204. Available at:
\url{https://doi.org/10.1007/s10664-020-09869-w}.

\bibitem[\citeproctext]{ref-ampatzoglouSDK4EDPlatformTechnical2022}
Ampatzoglou, A., Chatzigeorgiou, A., Arvanitou, E.M. and Bibi, S. (2022)
\textquotesingle{}{SDK4ED}: {A} platform for technical debt
management\textquotesingle, \emph{Software: Practice and Experience},
52(8), pp. 1879--1902. Available at:
\url{https://doi.org/10.1002/spe.3093}.

\bibitem[\citeproctext]{ref-avgeriouTechnicalDebtManagement2023}
Avgeriou, P., Ozkaya, I., Chatzigeorgiou, A., Ciolkowski, M., Ernst,
N.A., Koontz, R.J., Poort, E. and Shull, F. (2023)
\textquotesingle Technical {Debt Management}: {The Road Ahead} for
{Successful Software Delivery}\textquotesingle, \emph{2023 {IEEE}/{ACM
International Conference} on {Software Engineering}: {Future} of
{Software Engineering} ({ICSE-FoSE})}. \emph{2023 {IEEE}/{ACM
International Conference} on {Software Engineering}: {Future} of
{Software Engineering} ({ICSE-FoSE})}, Melbourne, Australia: IEEE, pp.
15--30. Available at:
\url{https://doi.org/10.1109/ICSE-FoSE59343.2023.00007}.

\bibitem[\citeproctext]{ref-bankerTechnicalDebtFirm2020}
Banker, R., Liang, Y. and Ramasubbu, N. (2020)
\textquotesingle Technical {Debt} and {Firm
Performance}\textquotesingle, \emph{Management Science}, 67(5), pp.
3174--3194. Available at: \url{https://doi.org/10.1287/mnsc.2019.3542}.

\bibitem[\citeproctext]{ref-bediUnderstandingFactorsAffecting2022}
Bedi, J. and Kaur, K. (2022) \textquotesingle Understanding factors
affecting technical debt\textquotesingle, \emph{International Journal of
Information Technology}, 14(2), pp. 1051--1060. Available at:
\url{https://doi.org/10.1007/s41870-020-00487-9}.

\bibitem[\citeproctext]{ref-beskerTechnicalDebtEmpirical2020}
Besker, T. (2020) \emph{Technical {Debt}: {An} empirical investigation
of its harmfulness and on management strategies in industry}. Chalmers
University of Technology. Available at:
\url{https://research.chalmers.se/publication/518318/file/518318_Fulltext.pdf}.

\bibitem[\citeproctext]{ref-biazottoTechnicalDebtManagement2024}
Biazotto, J.P., Feitosa, D., Avgeriou, P. and Nakagawa, E.Y. (2024)
\textquotesingle Technical debt management automation: {State} of the
art and future perspectives\textquotesingle, \emph{Information and
Software Technology}, 167, p. 107375. Available at:
\url{https://doi.org/10.1016/j.infsof.2023.107375}.

\bibitem[\citeproctext]{ref-ciardielloComparisonTOPSISSAW2023}
Ciardiello, F. and Genovese, A. (2023) \textquotesingle A comparison
between {TOPSIS} and {SAW} methods\textquotesingle, \emph{Annals of
Operations Research}, 325(2), pp. 967--994. Available at:
\url{https://doi.org/10.1007/s10479-023-05339-w}.

\bibitem[\citeproctext]{ref-CodeConductBCS2022}
\emph{Code of {Conduct} for {BCS Members}} (2022). Swindon, UK: BCS, The
Chartered Institute for IT, pp. 1--5. Available at:
\url{https://www.bcs.org/media/2211/bcs-code-of-conduct.pdf} (Accessed:
April 15, 2026).

\bibitem[\citeproctext]{ref-cunninghamWyCashPortfolioManagement1992}
Cunningham, W. (1992) \textquotesingle The {WyCash} portfolio management
system\textquotesingle, \emph{ACM SIGPLAN OOPS Messenger}, 4(2), pp.
29--30. Available at: \url{https://doi.org/10.1145/157710.157715}.

\bibitem[\citeproctext]{ref-FederalActData2023}
\emph{Federal {Act} on {Data Protection}} (2023). Available at:
\url{https://www.fedlex.admin.ch/eli/cc/2022/491/en} (Accessed: April
15, 2026).

\bibitem[\citeproctext]{ref-finmaCircular202312024}
\emph{{FINMA Circular} 2023/1 - {Operational} risks and resilience --
banks} (2024). Available at:
\url{https://www.finma.ch/en/~/media/finma/dokumente/dokumentencenter/myfinma/rundschreiben/finma-rs-2023-01-20221207.pdf}
(Accessed: March 29, 2026).

\bibitem[\citeproctext]{ref-forsgrenAccelerate2018}
Forsgren, N., Humble, J. and Kim, G. (2018) \emph{Accelerate: {The
Science} of {Lean Software} and {DevOps}: {Building} and {Scaling High
Performing Technology Organizations}}. Portland, OR: IT Revolution
Press.

\bibitem[\citeproctext]{ref-fowlerTechnicalDebtQuadrant2009}
Fowler, M. (2009) \emph{Technical {Debt Quadrant}}. Technical Debt
Quadrant. Available at:
\url{https://martinfowler.com/bliki/TechnicalDebtQuadrant.html}
(Accessed: March 29, 2026).

\bibitem[\citeproctext]{ref-fungprasertkulTechnicalDebtMonitoring2024}
Fungprasertkul, S., Bahsoon, R. and Kazman, R. (2024)
\textquotesingle Technical {Debt Monitoring Decision Making} with {Skin}
in the {Game}\textquotesingle, \emph{ACM Transactions on Software
Engineering and Methodology}, 33(7), pp. 1--27. Available at:
\url{https://doi.org/10.1145/3664805}.

\bibitem[\citeproctext]{ref-garcia-cascalesRankReversalTOPSIS2012}
García-Cascales, M.S. and Lamata, M.T. (2012) \textquotesingle On rank
reversal and {TOPSIS} method\textquotesingle, \emph{Mathematical and
Computer Modelling}, 56(5), pp. 123--132. Available at:
\url{https://doi.org/10.1016/j.mcm.2011.12.022}.

\bibitem[\citeproctext]{ref-grecoMethodologicalFrameworkComposite2019}
Greco, S., Ishizaka, A., Tasiou, M. and Torrisi, G. (2019)
\textquotesingle On the {Methodological Framework} of {Composite
Indices}: {A Review} of the {Issues} of {Weighting}, {Aggregation}, and
{Robustness}\textquotesingle, \emph{Social Indicators Research}, 141(1),
pp. 61--94. Available at:
\url{https://doi.org/10.1007/s11205-017-1832-9}.

\bibitem[\citeproctext]{ref-hakiDigitalNudgingTechnical2023}
Haki, K., Rieder, A., Buchmann, L. and W.Schneider, A. (2023)
\textquotesingle Digital nudging for technical debt management at
{Credit Suisse}\textquotesingle, \emph{European Journal of Information
Systems}, 32(1), pp. 64--80. Available at:
\url{https://doi.org/10.1080/0960085X.2022.2088413}.

\bibitem[\citeproctext]{ref-hanssenIdentifyingScalabilityDebt2019}
Hanssen, G.K., Brataas, G. and Martini, A. (2019)
\textquotesingle Identifying {Scalability Debt} in {Open
Systems}\textquotesingle, \emph{2019 {IEEE}/{ACM International
Conference} on {Technical Debt} ({TechDebt})}. \emph{2019 {IEEE}/{ACM
International Conference} on {Technical Debt} ({TechDebt})}, pp. 48--52.
Available at: \url{https://doi.org/10.1109/TechDebt.2019.00014}.

\bibitem[\citeproctext]{ref-hevnerDesignScienceInformation2004}
Hevner, A.R., March, S.T., Park, J. and Ram, S. (2004)
\textquotesingle Design {Science} in {Information Systems
Research}\textquotesingle, \emph{MIS Quarterly}, 28(1), pp. 75--105.
Available at: \url{https://doi.org/10.2307/25148625}.

\bibitem[\citeproctext]{ref-hwangMultipleAttributeDecision1981}
Hwang, C.-L. and Yoon, K. (1981) \emph{Multiple {Attribute Decision
Making}}. Edited by M. Beckmann and H.P. Künzi. Berlin, Heidelberg:
Springer (Lecture {Notes} in {Economics} and {Mathematical Systems}).
Available at: \url{https://doi.org/10.1007/978-3-642-48318-9}.

\bibitem[\citeproctext]{ref-ishakAnalysisFuzzyAHPTOPSIS2020}
Ishak, A. and Wanli (2020) \textquotesingle Analysis of {Fuzzy
AHP-TOPSIS Methods} in {Multi Criteria Decision Making}: {Literature
Review}\textquotesingle, \emph{IOP Conference Series: Materials Science
and Engineering}, 1003(1), p. 012147. Available at:
\url{https://doi.org/10.1088/1757-899X/1003/1/012147}.

\bibitem[\citeproctext]{ref-kellyThreePrinciplesPragmatism2020}
Kelly, L.M. and Cordeiro, M. (2020) \textquotesingle Three principles of
pragmatism for research on organizational processes\textquotesingle,
\emph{Methodological Innovations}, 13(2), p. 2059799120937242. Available
at: \url{https://doi.org/10.1177/2059799120937242}.

\bibitem[\citeproctext]{ref-khanIntegratedMethodologyRanking2022}
Khan, S. and Purohit, L. (2022) \textquotesingle An {Integrated
Methodology} of {Ranking Based} on {PROMETHEE-CRITIC} and {TOPSIS-CRITIC
In Web Service Domain}\textquotesingle, \emph{2022 {IEEE} 11th
{International Conference} on {Communication Systems} and {Network
Technologies} ({CSNT})}. \emph{2022 {IEEE} 11th {International
Conference} on {Communication Systems} and {Network Technologies}
({CSNT})}, pp. 335--340. Available at:
\url{https://doi.org/10.1109/CSNT54456.2022.9787620}.

\bibitem[\citeproctext]{ref-kleinwaksTechnicalDebtSystems2023}
Kleinwaks, H., Batchelor, A. and Bradley, T.H. (2023)
\textquotesingle Technical debt in systems engineering---{A} systematic
literature review\textquotesingle, \emph{Systems Engineering}, 26(5),
pp. 675--687. Available at: \url{https://doi.org/10.1002/sys.21681}.

\bibitem[\citeproctext]{ref-kruchtenTechnicalDebtMetaphor2012}
Kruchten, P., Nord, R.L. and Ozkaya, I. (2012)
\textquotesingle Technical {Debt}: {From Metaphor} to {Theory} and
{Practice}\textquotesingle, \emph{IEEE Software}, 29(6), pp. 18--21.
Available at: \url{https://doi.org/10.1109/MS.2012.167}.

\bibitem[\citeproctext]{ref-kulaReleasingFastSlow2019}
Kula, E., Rastogi, A., Huijgens, H., Deursen, A. van and Gousios, G.
(2019) \textquotesingle Releasing fast and slow: An exploratory case
study at {ING}\textquotesingle, \emph{Proceedings of the 2019 27th {ACM
Joint Meeting} on {European Software Engineering Conference} and
{Symposium} on the {Foundations} of {Software Engineering}}. New York,
NY, USA: Association for Computing Machinery ({ESEC}/{FSE} 2019), pp.
785--795. Available at: \url{https://doi.org/10.1145/3338906.3338978}.

\bibitem[\citeproctext]{ref-lenarduzziSystematicLiteratureReview2021}
Lenarduzzi, V., Besker, T., Taibi, D., Martini, A. and Arcelli Fontana,
F. (2021) \textquotesingle A systematic literature review on {Technical
Debt} prioritization: {Strategies}, processes, factors, and
tools\textquotesingle, \emph{Journal of Systems and Software}, 171, p.
110827. Available at: \url{https://doi.org/10.1016/j.jss.2020.110827}.

\bibitem[\citeproctext]{ref-markkanenPatchManagementPlanning2023}
Markkanen, V. and Frantti, T. (2023) \textquotesingle Patch management
planning - towards one-to-one policy\textquotesingle, \emph{2023 10th
{International Conference} on {Dependable Systems} and {Their
Applications} ({DSA})}. \emph{2023 10th {International Conference} on
{Dependable Systems} and {Their Applications} ({DSA})}, pp. 60--69.
Available at: \url{https://doi.org/10.1109/DSA59317.2023.00018}.

\bibitem[\citeproctext]{ref-mcconnellManagingTechnicalDebt2008}
McConnell, S. (2008) \textquotesingle Managing {Technical
Debt}\textquotesingle. Construx Software. Available at:
\url{https://www.construx.com/wp-content/uploads/2019/02/CxWhitePaper_TechnicalDebt.pdf}
(Accessed: March 21, 2026).

\bibitem[\citeproctext]{ref-nielsenITPortfolioManagement2022}
Nielsen, M.E. and Skaarup, S. (2022) \textquotesingle{}{IT Portfolio}
management as a framework for managing {Technical Debt}: {Theoretical}
framework applied on a case study\textquotesingle, \emph{14th
{International Conference} on {Theory} and {Practice} of {Electronic
Governance}}. \emph{{ICEGOV} 2021: 14th {International Conference} on
{Theory} and {Practice} of {Electronic Governance}}, New York, NY, USA:
Association for Computing Machinery, pp. 89--96. Available at:
\url{https://doi.org/10.1145/3494193.3494296}.

\bibitem[\citeproctext]{ref-nordExamplesTechnicalDebts2021}
Nord, R.L., Ozkaya, I. and Woody, C. (2021) \emph{Examples of {Technical
Debt}'s {Cybersecurity Impact}}. Pittsburgh, PA, USA: Carnegie Mellon
University, pp. 1--21. Available at:
\url{https://apps.dtic.mil/sti/trecms/pdf/AD1144728.pdf} (Accessed:
April 27, 2026).

\bibitem[\citeproctext]{ref-nursePatchNotPatch2025}
Nurse, J.R.C. (2025) \textquotesingle To {Patch} or~{Not} to~{Patch}:
{Motivations}, {Challenges}, and~{Implications}
for~{Cybersecurity}\textquotesingle, in A. Moallem (ed.) \emph{{HCI} for
{Cybersecurity}, {Privacy} and {Trust}}. Cham: Springer Nature
Switzerland, pp. 265--281. Available at:
\url{https://doi.org/10.1007/978-3-031-92833-8_16}.

\bibitem[\citeproctext]{ref-pandeyReviewTOPSISMethod2023}
Pandey, V., Komal and Dincer, H. (2023) \textquotesingle A review on
{TOPSIS} method and its extensions for different applications with
recent development\textquotesingle, \emph{Soft Computing}, 27(23), pp.
18011--18039. Available at:
\url{https://doi.org/10.1007/s00500-023-09011-0}.

\bibitem[\citeproctext]{ref-paudelExploringRelationshipTechnical2025}
Paudel, B., Gonzalez-Huerta, J., Zabardast, E. and Klotins, E. (2025)
\textquotesingle Exploring the {Relationship Between Technical Debt} and
{Lead Time}: {An Industrial Case Study}\textquotesingle, \emph{2025
{IEEE International Conference} on {Software Analysis}, {Evolution} and
{Reengineering} ({SANER})}. \emph{2025 {IEEE International Conference}
on {Software Analysis}, {Evolution} and {Reengineering} ({SANER})}, pp.
693--703. Available at:
\url{https://doi.org/10.1109/SANER64311.2025.00071}.

\bibitem[\citeproctext]{ref-pinaTechnicalDebtPrioritization2022}
Pina, D., Seaman, C. and Goldman, A. (2022) \textquotesingle Technical
{Debt Prioritization}: {A Developer}'s {Perspective}\textquotesingle,
\emph{2022 {IEEE}/{ACM International Conference} on {Technical Debt}
({TechDebt})}. \emph{2022 {IEEE}/{ACM International Conference} on
{Technical Debt} ({TechDebt})}, pp. 46--55. Available at:
\url{https://doi.org/10.1145/3524843.3528096}.

\bibitem[\citeproctext]{ref-ramasubbuControllingTechnicalDebt2021}
Ramasubbu, N. and Kemerer, C.F. (2021) \textquotesingle Controlling
{Technical Debt Remediation} in {Outsourced Enterprise Systems
Maintenance}: {An Empirical Analysis}\textquotesingle, \emph{Journal of
Management Information Systems}, 38(1), pp. 4--28. Available at:
\url{https://doi.org/10.1080/07421222.2021.1870377}.

\bibitem[\citeproctext]{ref-ramasubbuManagingTechnicalDebt2015}
Ramasubbu, N., Kemerer, C.F. and Woodard, C.J. (2015)
\textquotesingle Managing {Technical Debt}: {Insights} from {Recent
Empirical Evidence}\textquotesingle, \emph{IEEE Software}, 32(2), pp.
22--25. Available at: \url{https://doi.org/10.1109/MS.2015.45}.

\bibitem[\citeproctext]{ref-rinta-kahilaGettingTrappedTechnical2023}
Rinta-Kahila, T., Penttinen, E., Aalto University, Lyytinen, K. and Case
Western Reserve University (2023) \textquotesingle Getting {Trapped} in
{Technical Debt}: {Sociotechnical Analysis} of a {Legacy System}'s
{Replacement}\textquotesingle, \emph{MIS Quarterly}, 47(1), pp. 1--32.
Available at: \url{https://doi.org/10.25300/MISQ/2022/16711}.

\bibitem[\citeproctext]{ref-rollandExploringTensionsManagement2021}
Rolland, K.-H. and Lyytinen, K. (2021) \textquotesingle Exploring the
{Tensions Between Management} of {Architectural Debt} and {Digital
Innovation}: {The Case} of a {Financial Organization}\textquotesingle,
\emph{Proceedings of the 54th {Hawaii International Conference} on
{System Sciences}}. \emph{54th {Hawaii International Conference} on
{System Sciences}}, Hawaii International Conference on System Sciences,
pp. 6722--6732. Available at:
\url{https://doi.org/10.24251/HICSS.2021.807}.

\bibitem[\citeproctext]{ref-shekhovtsovComparativeCaseStudy2020}
Shekhovtsov, A. and Sałabun, W. (2020) \textquotesingle A comparative
case study of the {VIKOR} and {TOPSIS} rankings
similarity\textquotesingle, \emph{Procedia Computer Science}, 176, pp.
3730--3740. Available at:
\url{https://doi.org/10.1016/j.procs.2020.09.014}.

\bibitem[\citeproctext]{ref-stochelBusinessdrivenTechnicalDebt2023}
Stochel, M.G., Borek, T., Wawrowski, M.R. and Chołda, P. (2023)
\textquotesingle Business-driven technical debt management using
{Continuous Debt Valuation Approach} ({CoDVA})\textquotesingle,
\emph{Information and Software Technology}, 164, p. 107333. Available
at: \url{https://doi.org/10.1016/j.infsof.2023.107333}.

\bibitem[\citeproctext]{ref-tizioSoftwareUpdatesStrategies2023}
Tizio, G.D., Armellini, M. and Massacci, F. (2023)
\textquotesingle Software {Updates Strategies}: {A Quantitative
Evaluation Against Advanced Persistent Threats}\textquotesingle,
\emph{IEEE Transactions on Software Engineering}, 49(3), pp. 1359--1373.
Available at: \url{https://doi.org/10.1109/TSE.2022.3176674}.

\bibitem[\citeproctext]{ref-detoledoReducingIncidentsMicroservices2021}
de Toledo, S.S., Martini, A., Sjøberg, D.I.K., Przybyszewska, A. and
Frandsen, J.S. (2021) \textquotesingle Reducing {Incidents} in
{Microservices} by {Repaying Architectural Technical
Debt}\textquotesingle, \emph{2021 47th {Euromicro Conference} on
{Software Engineering} and {Advanced Applications} ({SEAA})}. \emph{2021
47th {Euromicro Conference} on {Software Engineering} and {Advanced
Applications} ({SEAA})}, pp. 196--205. Available at:
\url{https://doi.org/10.1109/SEAA53835.2021.00033}.

\bibitem[\citeproctext]{ref-tornhillCodeRedBusiness2022}
Tornhill, A. and Borg, M. (2022) \textquotesingle Code {Red}: {The
Business Impact} of {Code Quality} - {A Quantitative Study} of 39
{Proprietary Production Codebases}\textquotesingle, \emph{2022
{IEEE}/{ACM International Conference} on {Technical Debt} ({TechDebt})}.
\emph{2022 {IEEE}/{ACM International Conference} on {Technical Debt}
({TechDebt})}, pp. 11--20. Available at:
\url{https://doi.org/10.1145/3524843.3528091}.

\bibitem[\citeproctext]{ref-turnbullUseCaseStudy2021}
Turnbull, D., Chugh, R. and Luck, J. (2021) \textquotesingle The {Use}
of {Case Study Design} in {Learning Management System Research}: {A
Label} of {Convenience}?\textquotesingle, \emph{International Journal of
Qualitative Methods}, 20, pp. 1--11. Available at:
\url{https://doi.org/10.1177/16094069211004148}.

\bibitem[\citeproctext]{ref-vidoniInfiniteTechnicalDebt2022}
Vidoni, M., Codabux, Z. and Fard, F.H. (2022) \textquotesingle Infinite
technical debt\textquotesingle, \emph{Journal of Systems and Software},
190, p. 111336. Available at:
\url{https://doi.org/10.1016/j.jss.2022.111336}.

\bibitem[\citeproctext]{ref-vombrockeIntroductionDesignScience2020}
Vom Brocke, J., Hevner, A. and Maedche, A. (2020)
\textquotesingle Introduction to {Design Science
Research}\textquotesingle, in J. Vom Brocke, A. Hevner, and A. Maedche
(eds.) \emph{Design {Science Research}. {Cases}}. Cham: Springer
International Publishing, pp. 1--13. Available at:
\url{https://doi.org/10.1007/978-3-030-46781-4_1}.

\bibitem[\citeproctext]{ref-wieseITManagersPerspective2023}
Wiese, M. and Borowa, K. (2023) \textquotesingle{}{IT} managers'
perspective on {Technical Debt Management}\textquotesingle,
\emph{Journal of Systems and Software}, 202, p. 111700. Available at:
\url{https://doi.org/10.1016/j.jss.2023.111700}.

\bibitem[\citeproctext]{ref-yangIngeniousSolutionRank2020}
Yang, W. (2020) \textquotesingle Ingenious {Solution} for the {Rank
Reversal Problem} of {TOPSIS Method}\textquotesingle, \emph{Mathematical
Problems in Engineering}, 2020(1), p. 9676518. Available at:
\url{https://doi.org/10.1155/2020/9676518}.

\bibitem[\citeproctext]{ref-yangTaxonomyTechnicalDebt2019}
Yang, Y., Michel, R., Wade, J., Verma, D., Törngren, M. and Alelyani, T.
(2019) \textquotesingle Towards a taxonomy of technical debt for
{COTS-intensive} cyber physical systems\textquotesingle, \emph{Procedia
Computer Science}, 153, pp. 108--117. Available at:
\url{https://doi.org/10.1016/j.procs.2019.05.061}.

\bibitem[\citeproctext]{ref-zanakisMultiattributeDecisionMaking1998}
Zanakis, S.H., Solomon, A., Wishart, N. and Dublish, S. (1998)
\textquotesingle Multi-attribute decision making: {A} simulation
comparison of select methods\textquotesingle, \emph{European Journal of
Operational Research}, 107(3), pp. 507--529. Available at:
\url{https://doi.org/10.1016/S0377-2217(97)00147-1}.

\bibitem[\citeproctext]{ref-zytoonDecisionSupportModel2020}
Zytoon, M.A. (2020) \textquotesingle A {Decision Support Model} for
{Prioritization} of {Regulated Safety Inspections Using Integrated
Delphi}, {AHP} and {Double-Hierarchical TOPSIS
Approach}\textquotesingle, \emph{IEEE Access}, 8, pp. 83444--83464.
Available at: \url{https://doi.org/10.1109/ACCESS.2020.2991179}.

\end{CSLReferences}

\cleardoublepage
\phantomsection
\addcontentsline{toc}{part}{Appendices}
\appendix

\chapter{Appendix A}\label{appendix-a}

Lorem ipsum




\end{document}
